\section{Casi d'Uso}

\subsection{Attori}
\begin{itemize}
	\item \textbf{Utente}: Attore\textsubscript{G} primario $\rightarrow$ Rappresenta l'utilizzatore finale che interagisce con l'interfaccia dell'applicazione. 
    In base al capitolato e ai verbali esterni, il profilo di riferimento è un utente giovane, tipicamente uno studente, che utilizza l'applicazione per prendere appunti e desidera supporto nell'elaborazione tramite LLM (VE-5.3).
	\item \textbf{LLM (Large Language Model\textsubscript{G})}: Attore secondario (sistema esterno) $\rightarrow$ Fornisce i servizi di analisi e generazione testo tramite API compatibile con lo standard OpenAI. I modelli di riferimento sono Gemma, ChatGPT-OSS e Qwen, messi a disposizione dall'azienda proponente (VE-1.1).
\end{itemize}

% --------------------------------------------------------------------------
% 1. Macroarea 1: Gestione Editor e Markdown
% --------------------------------------------------------------------------
\newpage
\subsection{Macroarea 1 — Gestione Editor e Markdown}
\textit{Funzionalità core per la manipolazione e visualizzazione del testo.}

% UC1 - Scrittura di testo nell'editor
\hypertarget{uc1}{}
\subsubsection{UC 1: Scrittura testo nell'editor}

% use case diagram
\begin{figure}[H]
    \centering
    \includegraphics[width=0.5\textwidth]{scrittura_nota.png}
    \caption{UC 1 – Scrittura di testo nell'editor}
\end{figure}

\begin{itemize}
    \item \textbf{Attore Principale}: Utente
    \item \textbf{Precondizioni}: L'editor è attivo ed è presente una nota aperta nell'editor.
    \item \textbf{Scenario Principale\textsubscript{G}}:
    \begin{enumerate}
        \item L'utente posiziona il cursore nell'area di testo.
        \item L'utente inserisce testo tramite tastiera.
        \item Il sistema aggiorna in tempo reale l'area di testo con il contenuto inserito.
        \item Il sistema aggiorna l'anteprima renderizzata, se la modalità di visualizzazione attiva lo prevede.
    \end{enumerate}
    \item \textbf{Postcondizioni}: Il testo inserito è visibile nell'editor e, ove applicabile, nell'anteprima renderizzata.
    \item \textbf{Trigger}: L'utente vuole inserire o modificare i contenuti testuali della nota.
\end{itemize}

% UC2 - Copia testo
\hypertarget{uc2}{}
\subsubsection{UC 2: Copia testo}

% use case diagram
\begin{figure}[H]
    \centering
    \includegraphics[width=0.6\textwidth]{copia_testo.png}
    \caption{UC 2 – Copia testo}
\end{figure}

\begin{itemize}
    \item \textbf{Attore Principale}: Utente
    \item \textbf{Precondizioni}: L'editor è attivo, è presente una nota aperta nell'editor ed una porzione di testo è selezionata.
    \item \textbf{Scenario Principale}:
        \begin{enumerate}
        \item L'utente attiva il comando di copia.
        \item Il sistema copia il testo selezionato negli appunti\textsubscript{G}.
        \end{enumerate}
    \item \textbf{Postcondizioni}: Il testo è stato copiato negli appunti.
    \item \textbf{Trigger}: L'utente vuole copiare una porzione di testo.
\end{itemize}

% UC3 - Taglia testo
\hypertarget{uc3}{}
\subsubsection{UC 3: Taglia testo}

% use case diagram
\begin{figure}[H]
    \centering
    \includegraphics[width=0.6\textwidth]{taglia_testo.png}
    \caption{UC 3 – Taglia testo}
\end{figure}

\begin{itemize}
    \item \textbf{Attore Principale}: Utente
    \item \textbf{Precondizioni}: L'editor è attivo, è presente una nota aperta nell'editor ed una porzione di testo è selezionata.
    \item \textbf{Scenario Principale}:
    \begin{enumerate}
        \item L'utente attiva il comando di taglio.
        \item Il sistema rimuove il testo selezionato dall'editor e lo inserisce negli appunti.
    \end{enumerate}
    \item \textbf{Postcondizioni}: Il testo è stato rimosso dall'editor e copiato negli appunti.
    \item \textbf{Trigger}: L'utente vuole spostare una porzione di testo.
\end{itemize}

% UC4 - Incolla testo
\hypertarget{uc4}{}
\subsubsection{UC 4: Incolla testo}

% use case diagram
\begin{figure}[H]
    \centering
    \includegraphics[width=0.6\textwidth]{incolla_testo.png}
    \caption{UC 4 – Incolla testo}
\end{figure}

\begin{itemize}
    \item \textbf{Attore Principale}: Utente
    \item \textbf{Precondizioni}: L'editor è attivo, è presente una nota aperta nell'editor ed una porzione di testo è stata copiata o tagliata negli appunti.
    \item \textbf{Scenario Principale}:
    \begin{enumerate}
        \item L'utente posiziona il cursore nel punto desiderato.
        \item L'utente attiva il comando di incolla.
        \item Il sistema inserisce il contenuto degli appunti nella posizione del cursore.
    \end{enumerate}
    \item \textbf{Postcondizioni}: Il testo è stato incollato nella posizione scelta.
    \item \textbf{Trigger}: L'utente vuole inserire testo precedentemente copiato o tagliato.
\end{itemize}

%---------------GRASSETTO----------------

% UC5 - Attivazione grassetto 
\hypertarget{uc5}{}
\subsubsection{UC 5: Attivazione grassetto}

% use case diagram
\begin{figure}[H]
    \centering
    \includegraphics[width=0.6\textwidth]{attivazione_grassetto.png}
    \caption{UC 5 – Attivazione grassetto }
\end{figure}

\begin{itemize}
    \item \textbf{Attore Principale}: Utente
    \item \textbf{Precondizioni}: L'editor è attivo ed è presente una nota aperta nell'editor.
    \item \textbf{Scenario Principale}:
        \begin{enumerate}
        \item L'utente attiva il comando "Grassetto".
        \item Il sistema abilita la modalità di inserimento in grassetto.
        \end{enumerate}
    \item \textbf{Postcondizioni}: Il testo inserito successivamente viene formattato in grassetto.
    \item \textbf{Trigger}: L'utente desidera applicare la formattazione grassetto durante l'inserimento del testo.
    \end{itemize}


% UC6 - Applicazione grassetto al testo selezionato
\hypertarget{uc6}{}
\subsubsection{UC 6: Applicazione grassetto al testo selezionato}

% use case diagram
\begin{figure}[H]
    \centering
    \includegraphics[width=0.5\textwidth]{applicazione_grassetto.png}
    \caption{UC 6 – Applicazione grassetto al testo selezionato}
\end{figure}

\begin{itemize}
    \item \textbf{Attore Principale}: Utente
    \item \textbf{Precondizioni}: L'editor è attivo, è presente una nota aperta nell'editor ed una porzione di testo è selezionata.
    \item \textbf{Scenario Principale}:
        \begin{enumerate}
        \item L'utente attiva il comando "Grassetto".
        \item Il sistema applica la formattazione grassetto al testo selezionato.
        \end{enumerate}
    \item \textbf{Postcondizioni}: La porzione di testo selezionata presenta la formattazione grassetto.
    \item \textbf{Trigger}: L'utente desidera applicare il formato grassetto a testo esistente.
\end{itemize}

% UC7 -  Disattivazione grassetto 
\hypertarget{uc7}{}
\subsubsection{UC 7: Disattivazione grassetto }

% use case diagram
\begin{figure}[H]
    \centering
    \includegraphics[width=0.5\textwidth]{disattivazione_grassetto.png}
    \caption{UC 7 – Disattivazione grassetto }
\end{figure}

\begin{itemize}
\item \textbf{Attore Principale}: Utente
\item \textbf{Precondizioni}: La modalità di inserimento in grassetto è attiva.
\item \textbf{Scenario Principale}:
    \begin{enumerate}
    \item L'utente disattiva il comando "Grassetto".
    \item Il sistema disabilita la modalità di inserimento in grassetto.
    \end{enumerate}
\item \textbf{Postcondizioni}: Il testo inserito successivamente non verrà formattato in grassetto.
\item \textbf{Trigger}: L'utente desidera interrompere la scrittura in grassetto.
\end{itemize}

% UC8 - Rimozione grassetto dal testo selezionato
\hypertarget{uc8}{}
\subsubsection{UC 8: Rimozione grassetto dal testo selezionato}

% use case diagram
\begin{figure}[H]
    \centering
    \includegraphics[width=0.5\textwidth]{rimozione_grassetto.png}
    \caption{UC 8 – Rimozione grassetto dal testo selezionato}
\end{figure}

\begin{itemize}
    \item \textbf{Attore Principale}: Utente
    \item \textbf{Precondizioni}: L'editor è attivo, è presente una nota aperta nell'editor ed una porzione di testo formattata in grassetto è selezionata.
    \item \textbf{Scenario Principale}:
    \begin{enumerate}
        \item L'utente disattiva il comando "Grassetto".
        \item Il sistema rimuove la formattazione grassetto dal testo selezionato.
    \end{enumerate}
    \item \textbf{Postcondizioni}: La porzione di testo selezionata non presenta più la formattazione grassetto.
    \item \textbf{Trigger}: L'utente desidera rimuovere il formato grassetto da testo esistente.
\end{itemize}

%---------------CORSIVO----------------

% UC9 -  Attivazione corsivo 
\hypertarget{uc9}{}
\subsubsection{UC 9: Attivazione corsivo }

% use case diagram
\begin{figure}[H]
    \centering
    \includegraphics[width=0.5\textwidth]{attivazione_corsivo.png}
    \caption{UC 9 – Attivazione corsivo }
\end{figure}

\begin{itemize}
    \item \textbf{Attore Principale}: Utente
    \item \textbf{Precondizioni}: L'editor è attivo ed è presente una nota aperta nell'editor.
    \item \textbf{Scenario Principale}:
        \begin{enumerate}
        \item L'utente attiva il comando "Corsivo".
        \item Il sistema abilita la modalità di inserimento  in corsivo.
        \end{enumerate}
    \item \textbf{Postcondizioni}: Il testo inserito successivamente viene formattato in corsivo.
    \item \textbf{Trigger}: L'utente desidera applicare la formattazione corsivo durante l'inserimento del testo.
    \end{itemize}

% UC10 - Applicazione corsivo al testo selezionato
\hypertarget{uc10}{}
\subsubsection{UC 10: Applicazione corsivo al testo selezionato}

% use case diagram
\begin{figure}[H]
    \centering
    \includegraphics[width=0.5\textwidth]{applicazione_corsivo.png}
    \caption{UC 10 – Applicazione corsivo al testo selezionato}
\end{figure}

\begin{itemize}
    \item \textbf{Attore Principale}: Utente
    \item \textbf{Precondizioni}: L'editor è attivo, è presente una nota aperta nell'editor ed una porzione di testo è selezionata.
    \item \textbf{Scenario Principale}:
        \begin{enumerate}
        \item L'utente attiva il comando "Corsivo".
        \item Il sistema applica la formattazione corsivo al testo selezionato.
        \end{enumerate}
    \item \textbf{Postcondizioni}: La porzione di testo selezionata presenta la formattazione corsivo.
    \item \textbf{Trigger}: L'utente desidera applicare il formato corsivo a testo esistente.
\end{itemize}

% UC11 -  Disattivazione corsivo 
\hypertarget{uc11}{}
\subsubsection{UC 11: Disattivazione corsivo }

% use case diagram
\begin{figure}[H]
    \centering
    \includegraphics[width=0.5\textwidth]{disattivazione_corsivo.png}
    \caption{UC 11 – Disattivazione corsivo }
\end{figure}

\begin{itemize}
\item \textbf{Attore Principale}: Utente
\item \textbf{Precondizioni}: La modalità di inserimento in corsivo è attiva.
\item \textbf{Scenario Principale}:
    \begin{enumerate}
    \item L'utente disattiva il comando "Corsivo".
    \item Il sistema disabilita la modalità di inserimento in corsivo.
    \end{enumerate}
\item \textbf{Postcondizioni}: Il testo inserito successivamente non verrà formattato in corsivo.
\item \textbf{Trigger}: L'utente desidera interrompere la scrittura in corsivo.
\end{itemize}

% UC12 - Rimozione corsivo dal testo selezionato
\hypertarget{uc12}{}
\subsubsection{UC 12: Rimozione corsivo dal testo selezionato}

% use case diagram
\begin{figure}[H]
    \centering
    \includegraphics[width=0.5\textwidth]{rimozione_corsivo.png}
    \caption{UC 12 – Rimozione corsivo dal testo selezionato}
\end{figure}

\begin{itemize}
    \item \textbf{Attore Principale}: Utente
    \item \textbf{Precondizioni}: L'editor è attivo, è presente una nota aperta nell'editor ed una porzione di testo formattata in corsivo è selezionata.
    \item \textbf{Scenario Principale}:
    \begin{enumerate}
        \item L'utente disattiva il comando "Corsivo".
        \item Il sistema rimuove la formattazione corsivo dal testo selezionato.
    \end{enumerate}
    \item \textbf{Postcondizioni}: La porzione di testo selezionata non presenta più la formattazione corsivo.
    \item \textbf{Trigger}: L'utente desidera rimuovere il formato corsivo da testo esistente.
\end{itemize}

%---------------SOTTOLINEATO----------------

% UC13 -  Attivazione sottolineato 
\hypertarget{uc13}{}
\subsubsection{UC 13: Attivazione sottolineato }

% use case diagram
\begin{figure}[H]
    \centering
    \includegraphics[width=0.5\textwidth]{attivazione_sottolineato.png}
    \caption{UC 13 – Attivazione sottolineato }
\end{figure}

\begin{itemize}
    \item \textbf{Attore Principale}: Utente
    \item \textbf{Precondizioni}: L'editor è attivo ed è presente una nota aperta nell'editor.
    \item \textbf{Scenario Principale}:
        \begin{enumerate}
        \item L'utente attiva il comando "Sottolineato".
        \item Il sistema abilita la modalità di inserimento  in sottolineato.
        \end{enumerate}
    \item \textbf{Postcondizioni}: Il testo inserito successivamente viene formattato in sottolineato.
    \item \textbf{Trigger}: L'utente desidera applicare la formattazione sottolineato durante l'inserimento del testo.
    \end{itemize}

% UC14 - Applicazione sottolineato al testo selezionato
\hypertarget{uc14}{}
\subsubsection{UC 14: Applicazione sottolineato al testo selezionato}

% use case diagram
\begin{figure}[H]
    \centering
    \includegraphics[width=0.5\textwidth]{applicazione_sottolineato.png}
    \caption{UC 14 – Applicazione sottolineato al testo selezionato}
\end{figure}

\begin{itemize}
    \item \textbf{Attore Principale}: Utente
    \item \textbf{Precondizioni}: L'editor è attivo, è presente una nota aperta nell'editor ed una porzione di testo è selezionata.
    \item \textbf{Scenario Principale}:
        \begin{enumerate}
        \item L'utente attiva il comando "Sottolineato".
        \item Il sistema applica la formattazione sottolineato al testo selezionato.
        \end{enumerate}
    \item \textbf{Postcondizioni}: La porzione di testo selezionata presenta la formattazione sottolineato.
    \item \textbf{Trigger}: L'utente desidera applicare il formato sottolineato a testo esistente.
\end{itemize}

% UC15 -  Disattivazione sottolineato 
\hypertarget{uc15}{}
\subsubsection{UC 15: Disattivazione sottolineato}

% use case diagram
\begin{figure}[H]
    \centering
    \includegraphics[width=0.5\textwidth]{disattivazione_sottolineato.png}
    \caption{UC 15 – Disattivazione sottolineato}
\end{figure}

\begin{itemize}
    \item \textbf{Attore Principale}: Utente
    \item \textbf{Precondizioni}: La modalità di inserimento in sottolineato è attiva.
    \item \textbf{Scenario Principale}:
        \begin{enumerate}
        \item L'utente disattiva il comando "Sottolineato".
        \item Il sistema disabilita la modalità di inserimento in sottolineato.
        \end{enumerate}
    \item \textbf{Postcondizioni}: Il testo inserito successivamente non verrà formattato in sottolineato.
    \item \textbf{Trigger}: L'utente desidera interrompere la scrittura in sottolineato.
\end{itemize}

% UC16 - Rimozione sottolineato dal testo selezionato
\hypertarget{uc16}{}
\subsubsection{UC 16: Rimozione sottolineato dal testo selezionato}

% use case diagram
\begin{figure}[H]
    \centering
    \includegraphics[width=0.5\textwidth]{rimozione_sottolineato.png}
    \caption{UC 16 – Rimozione sottolineato dal testo selezionato}
\end{figure}

\begin{itemize}
    \item \textbf{Attore Principale}: Utente
    \item \textbf{Precondizioni}: L'editor è attivo, è presente una nota aperta nell'editor ed una porzione di testo formattata in sottolineato è selezionata.
    \item \textbf{Scenario Principale}:
    \begin{enumerate}
        \item L'utente disattiva il comando "Sottolineato".
        \item Il sistema rimuove la formattazione sottolineato dal testo selezionato.
    \end{enumerate}
    \item \textbf{Postcondizioni}: La porzione di testo selezionata non presenta più la formattazione sottolineato.
    \item \textbf{Trigger}: L'utente desidera rimuovere il formato sottolineato da testo esistente.
\end{itemize}

%---------------BARRATO----------------

% UC17 -  Attivazione barrato
\hypertarget{uc17}{}
\subsubsection{UC 17: Attivazione barrato}

% use case diagram
\begin{figure}[H]
    \centering
    \includegraphics[width=0.5\textwidth]{attivazione_barrato.png}
    \caption{UC 17 – Attivazione barrato}
\end{figure}

\begin{itemize}
    \item \textbf{Attore Principale}: Utente
    \item \textbf{Precondizioni}: L'editor è attivo ed è presente una nota aperta nell'editor.
    \item \textbf{Scenario Principale}:
        \begin{enumerate}
        \item L'utente attiva il comando "Barrato".
        \item Il sistema abilita la modalità di inserimento in barrato.
        \end{enumerate}
    \item \textbf{Postcondizioni}: Il testo inserito successivamente viene formattato in barrato.
    \item \textbf{Trigger}: L'utente desidera applicare la formattazione barrato durante l'inserimento del testo.
\end{itemize}

% UC18 - Applicazione barrato al testo selezionato
\hypertarget{uc18}{}
\subsubsection{UC 18: Applicazione barrato al testo selezionato}

% use case diagram
\begin{figure}[H]
    \centering
    \includegraphics[width=0.5\textwidth]{applicazione_barrato.png}
    \caption{UC 18 – Applicazione barrato al testo selezionato}
\end{figure}

\begin{itemize}
    \item \textbf{Attore Principale}: Utente
    \item \textbf{Precondizioni}: L'editor è attivo, è presente una nota aperta nell'editor ed una porzione di testo è selezionata.
    \item \textbf{Scenario Principale}:
        \begin{enumerate}
        \item L'utente attiva il comando "Barrato".
        \item Il sistema applica la formattazione barrato al testo selezionato.
        \end{enumerate}
    \item \textbf{Postcondizioni}: La porzione di testo selezionata presenta la formattazione barrato.
    \item \textbf{Trigger}: L'utente desidera applicare il formato barrato a testo esistente.
\end{itemize}

% UC19 -  Disattivazione barrato
\hypertarget{uc19}{}
\subsubsection{UC 19: Disattivazione barrato}

% use case diagram
\begin{figure}[H]
    \centering
    \includegraphics[width=0.5\textwidth]{disattivazione_barrato.png}
    \caption{UC 19 – Disattivazione barrato}
\end{figure}

\begin{itemize}
    \item \textbf{Attore Principale}: Utente
    \item \textbf{Precondizioni}: La modalità di inserimento in barrato è attiva.
    \item \textbf{Scenario Principale}:
        \begin{enumerate}
        \item L'utente disattiva il comando "Barrato".
        \item Il sistema disabilita la modalità di inserimento in barrato.
        \end{enumerate}
    \item \textbf{Postcondizioni}: Il testo inserito successivamente non verrà formattato in barrato.
    \item \textbf{Trigger}: L'utente desidera interrompere la scrittura in barrato.
\end{itemize}

% UC20 - Rimozione barrato dal testo selezionato
\hypertarget{uc20}{}
\subsubsection{UC 20: Rimozione barrato dal testo selezionato}

% use case diagram
\begin{figure}[H]
    \centering
    \includegraphics[width=0.5\textwidth]{rimozione_barrato.png}
    \caption{UC 20 – Rimozione barrato dal testo selezionato}
\end{figure}

\begin{itemize}
    \item \textbf{Attore Principale}: Utente
    \item \textbf{Precondizioni}: L'editor è attivo, è presente una nota aperta nell'editor ed una porzione di testo formattata in barrato è selezionata.
    \item \textbf{Scenario Principale}:
    \begin{enumerate}
        \item L'utente disattiva il comando "Barrato".
        \item Il sistema rimuove la formattazione barrato dal testo selezionato.
    \end{enumerate}
    \item \textbf{Postcondizioni}: La porzione di testo selezionata non presenta più la formattazione barrato.
    \item \textbf{Trigger}: L'utente desidera rimuovere il formato barrato da testo esistente.
\end{itemize}

%---------------CITAZIONE----------------

% UC21 -  Attivazione formattazione citazione
\hypertarget{uc21}{}
\subsubsection{UC 21: Attivazione formattazione citazione}

% use case diagram
\begin{figure}[H]
    \centering
    \includegraphics[width=0.5\textwidth]{attivazione_citazione.png}
    \caption{UC 21 – Attivazione formattazione citazione}
\end{figure}

\begin{itemize}
    \item \textbf{Attore Principale}: Utente
    \item \textbf{Precondizioni}: L'editor è attivo ed è presente una nota aperta nell'editor.
    \item \textbf{Scenario Principale}:
        \begin{enumerate}
        \item L'utente attiva il comando "Citazione".
        \item Il sistema abilita la modalità di inserimento citazione.
        \end{enumerate}
    \item \textbf{Postcondizioni}: Il testo inserito successivamente viene formattato come citazione.
    \item \textbf{Trigger}: L'utente desidera scrivere nuovo testo come citazione.
\end{itemize}

% UC22 - Applicazione citazione al testo selezionato
\hypertarget{uc22}{}
\subsubsection{UC 22: Applicazione formattazione citazione al testo selezionato}

% use case diagram
\begin{figure}[H]
    \centering
    \includegraphics[width=0.5\textwidth]{applicazione_citazione.png}
    \caption{UC 22 – Applicazione formattazione citazione al testo selezionato}
\end{figure}

\begin{itemize}
    \item \textbf{Attore Principale}: Utente
    \item \textbf{Precondizioni}: L'editor è attivo, è presente una nota aperta nell'editor ed una porzione di testo è selezionata.
    \item \textbf{Scenario Principale}:
        \begin{enumerate}
        \item L'utente attiva il comando "Citazione".
        \item Il sistema applica la formattazione citazione al testo selezionato.
        \end{enumerate}
    \item \textbf{Postcondizioni}: La porzione di testo selezionata è convertita in blocco citazione.
    \item \textbf{Trigger}: L'utente desidera applicare la fomattazione citazione a testo esistente.
\end{itemize}

% UC23 -  Disattivazione citazione
\hypertarget{uc23}{}
\subsubsection{UC 23: Disattivazione citazione}

% use case diagram
\begin{figure}[H]
    \centering
    \includegraphics[width=0.5\textwidth]{disattivazione_citazione.png}
    \caption{UC 23 – Disattivazione citazione}
\end{figure}

\begin{itemize}
    \item \textbf{Attore Principale}: Utente
    \item \textbf{Precondizioni}: La modalità di inserimento citazione è attiva.
    \item \textbf{Scenario Principale}:
        \begin{enumerate}
        \item L'utente disattiva il comando "Citazione".
        \item Il sistema disabilita la modalità di inserimento citazione.
        \end{enumerate}
    \item \textbf{Postcondizioni}: Il testo inserito successivamente non verrà formattato come citazione.
    \item \textbf{Trigger}: L'utente desidera interrompere la modalità citazione.
\end{itemize}

% UC24 - Rimozione formattazione citazione dal testo selezionato
\hypertarget{uc24}{}
\subsubsection{UC 24: Rimozione formattazione citazione dal testo selezionato}

% use case diagram
\begin{figure}[H]
    \centering
    \includegraphics[width=0.5\textwidth]{rimozione_citazione.png}
    \caption{UC 24 – Rimozione formattazione citazione dal testo selezionato}
\end{figure}

\begin{itemize}
    \item \textbf{Attore Principale}: Utente
    \item \textbf{Precondizioni}: L'editor è attivo, è presente una nota aperta nell'editor ed una porzione di testo formattata come citazione è selezionata.
    \item \textbf{Scenario Principale}:
    \begin{enumerate}
        \item L'utente disattiva il comando "Citazione".
        \item Il sistema rimuove la formattazione citazione dal testo selezionato.
    \end{enumerate}
    \item \textbf{Postcondizioni}: La porzione di testo selezionata non presenta più la formattazione citazione.
    \item \textbf{Trigger}: L'utente desidera rimuovere la formattazione citazione da testo esistente.
\end{itemize}

%---------------TITOLO----------------

% UC25 - Inserimento titolo
\hypertarget{uc25}{}
\subsubsection{UC 25: Inserimento titolo}

% use case diagram
\begin{figure}[H]
    \centering
    \includegraphics[width=0.5\textwidth]{inserimento_titolo.png}
    \caption{UC 25 – Inserimento titolo}
\end{figure}

\begin{itemize}
    \item \textbf{Attore Principale}: Utente
    \item \textbf{Precondizioni}: L'editor è attivo ed è presente una nota aperta nell'editor.
    \item \textbf{Scenario Principale}:
        \begin{enumerate}
        \item L'utente attiva il comando "Titolo".
        \item Il sistema abilita la modalità titolo.
        \item L'utente inserisce il testo.
        \end{enumerate}    
    \item \textbf{Postcondizioni}: Il testo inserito viene formattato come titolo.
    \item \textbf{Trigger}: L'utente desidera creare una nuova intestazione nel documento.
\end{itemize}

% UC 26 - Applicazione formattazione titolo a testo selezionato
\hypertarget{uc26}{}
\subsubsection{UC 26: Applicazione formattazione titolo a testo selezionato}

% use case diagram
\begin{figure}[H]
    \centering
    \includegraphics[width=0.5\textwidth]{applicazione_titolo.png}
    \caption{UC 26 – Applicazione formattazione titolo a testo selezionato}
\end{figure}

\begin{itemize}
    \item \textbf{Attore Principale}: Utente
    \item \textbf{Precondizioni}: L'editor è attivo, è presente una nota aperta nell'editor ed una porzione di testo è selezionata.
    \item \textbf{Scenario Principale}:
        \begin{enumerate}
        \item L'utente attiva il comando "Titolo".
        \item Il sistema applica la formattazione titolo al testo selezionato.
        \end{enumerate}
    \item \textbf{Postcondizioni}: La porzione di testo selezionata presenta la formattazione titolo.
    \item \textbf{Trigger}: L'utente desidera applicare la formattazione titolo a testo esistente.
\end{itemize}

% UC 27 - Disattivazione titolo
\hypertarget{uc27}{}
\subsubsection{UC 27: Disattivazione titolo}

% use case diagram
\begin{figure}[H]
    \centering
    \includegraphics[width=0.5\textwidth]{disattivazione_titolo.png}
    \caption{UC 27 – Disattivazione titolo}
\end{figure}

\begin{itemize}
    \item \textbf{Attore Principale}: Utente
    \item \textbf{Precondizioni}: La modalità di inserimento titolo è attiva.
    \item \textbf{Scenario Principale}:
        \begin{enumerate}
        \item L'utente disattiva il comando "Titolo".
        \item Il sistema disabilita la modalità di inserimento titolo.
        \end{enumerate}
    \item \textbf{Postcondizioni}: Il testo inserito successivamente non verrà formattato come titolo.
    \item \textbf{Trigger}: L'utente desidera interrompere la modalità titolo.
\end{itemize}

% UC 28 - Rimozione formattazione titolo da testo selezionato
\hypertarget{uc28}{}
\subsubsection{UC 28: Rimozione formattazione titolo da testo selezionato}

% use case diagram
\begin{figure}[H]
    \centering
    \includegraphics[width=0.5\textwidth]{rimozione_titolo.png}
    \caption{UC 28 – Rimozione formattazione titolo su testo selezionato}
\end{figure}

\begin{itemize}
    \item \textbf{Attore Principale}: Utente
    \item \textbf{Precondizioni}: L'editor è attivo, è presente una nota aperta nell'editor ed una porzione di testo formattata come titolo è selezionata.
    \item \textbf{Scenario Principale}:
        \begin{enumerate}
        \item L'utente disattiva il comando "Titolo".
        \item Il sistema rimuove la formattazione titolo dal testo selezionato.
        \end{enumerate}
    \item \textbf{Postcondizioni}: La porzione di testo selezionata non presenta più la formattazione titolo.
    \item \textbf{Trigger}: L'utente desidera rimuovere la formattazione titolo da testo esistente.
\end{itemize}

%---------------LINK IPERTESTUALE----------------

% UC 29 - Inserimento link ipertestuale
\hypertarget{uc29}{}
\subsubsection{UC 29: Inserimento link ipertestuale}

% use case diagram
\begin{figure}[H]
    \centering
    \includegraphics[width=0.5\textwidth]{inserimento_link.png}
    \caption{UC 29 – Inserimento link ipertestuale}
\end{figure}

\begin{itemize}
    \item \textbf{Attore Principale}: Utente
    \item \textbf{Precondizioni}: L'editor è attivo ed è presente una nota aperta nell'editor.
    \item \textbf{Scenario Principale}:
    \begin{enumerate}
        \item L'utente attiva il comando "Inserisci link".
        \item Il sistema richiede l'inserimento dell'URL e del testo associato.
        \item L'utente inserisce i dati richiesti.
        \item Il sistema crea il collegamento.
    \end{enumerate}
    \item \textbf{Postcondizioni}: Il collegamento ipertestuale è inserito nel documento.
    \item \textbf{Trigger}: L'utente desidera inserire un collegamento nel documento.
\end{itemize}

% UC 30 - Modifica link
\hypertarget{uc30}{}
\subsubsection{UC 30: Modifica link}

% use case diagram
\begin{figure}[H]
    \centering
    \includegraphics[width=0.5\textwidth]{modifica_link.png}
    \caption{UC 30 – Modifica link}
\end{figure}

\begin{itemize}
    \item \textbf{Attore Principale}: Utente
    \item \textbf{Precondizioni}: L'editor è attivo, è presente una nota aperta nell'editor ed è presente almeno un collegamento nel documento.
    \item \textbf{Scenario Principale}:
    \begin{enumerate}
        \item L'utente seleziona un collegamento esistente.
        \item L'utente modifica i parametri del collegamento.
        \item Il sistema aggiorna il collegamento.
    \end{enumerate}
    \item \textbf{Postcondizioni}: Il collegamento presenta i nuovi parametri inseriti.
    \item \textbf{Trigger}: L'utente desidera modificare un collegamento esistente.
\end{itemize}

% UC 31 - Rimozione link
\hypertarget{uc31}{}
\subsubsection{UC 31: Rimozione link}

% use case diagram
\begin{figure}[H]
    \centering
    \includegraphics[width=0.5\textwidth]{rimozione_link.png}
    \caption{UC 31 – Rimozione link}
\end{figure}

\begin{itemize}
    \item \textbf{Attore Principale}: Utente
    \item \textbf{Precondizioni}: Un collegamento è presente nel documento ed è selezionato.
    \item \textbf{Scenario Principale}:
    \begin{enumerate}
        \item L'utente seleziona un collegamento.
        \item L'utente attiva il comando di rimozione.
        \item Il sistema elimina il collegamento.
    \end{enumerate}
    \item \textbf{Postcondizioni}: Il collegamento viene rimosso dal documento.
    \item \textbf{Trigger}: L'utente desidera rimuovere un collegamento esistente.
\end{itemize}

%---------------ELENCHI----------------

% UC 32 - Inserimento elenco puntato
\hypertarget{uc32}{}
\subsubsection{UC 32: Inserimento elenco puntato}

% use case diagram
\begin{figure}[H]
    \centering
    \includegraphics[width=0.5\textwidth]{inserimento_elenco_puntato.png}
    \caption{UC 32 – Inserimento elenco puntato}
\end{figure}

\begin{itemize}
    \item \textbf{Attore Principale}: Utente
    \item \textbf{Precondizioni}: L'editor è attivo ed è presente una nota aperta nell'editor.
    \item \textbf{Scenario Principale}:
    \begin{enumerate}
        \item L'utente attiva il comando "Elenco puntato".
        \item Il sistema abilita la modalità di inserimento elenco puntato.
        \item L'utente inserisce il contenuto del primo elemento.
    \end{enumerate}
    \item \textbf{Postcondizioni}: Il testo inserito viene formattato come primo elemento di un elenco puntato.
    \item \textbf{Trigger}: L'utente desidera creare un elenco puntato nel documento.
\end{itemize}

% UC 33 - Inserimento elenco numerato
\hypertarget{uc33}{}
\subsubsection{UC 33: Inserimento elenco numerato}

% use case diagram
\begin{figure}[H]
    \centering
    \includegraphics[width=0.5\textwidth]{inserimento_elenco_numerato.png}
    \caption{UC 33 – Inserimento elenco numerato}
\end{figure}

\begin{itemize}
    \item \textbf{Attore Principale}: Utente
    \item \textbf{Precondizioni}: L'editor è attivo ed è presente una nota aperta nell'editor.
    \item \textbf{Scenario Principale}:
    \begin{enumerate}
        \item L'utente attiva il comando "Elenco numerato".
        \item Il sistema abilita la modalità di inserimento elenco numerato.
        \item L'utente inserisce il contenuto del primo elemento.
    \end{enumerate}
    \item \textbf{Postcondizioni}: Il testo inserito viene formattato come primo elemento (numero 1) di un elenco numerato.
    \item \textbf{Trigger}: L'utente desidera creare un elenco numerato nel documento.
\end{itemize}

% UC 34 - Aggiunta elemento elenco
\hypertarget{uc34}{}
\subsubsection{UC 34: Aggiunta elemento elenco}

% use case diagram
\begin{figure}[H]
    \centering
    \includegraphics[width=0.5\textwidth]{aggiunta_elemento_elenco.png}
    \caption{UC 34 – Aggiunta elemento elenco}
\end{figure}

\begin{itemize}
    \item \textbf{Attore Principale}: Utente
    \item \textbf{Precondizioni}: L'utente sta scrivendo all'interno di un elenco (puntato o numerato) attivo.
    \item \textbf{Scenario Principale}:
    \begin{enumerate}
        \item L'utente richiede l'inserimento di un nuovo blocco di testo (es. premendo Invio).
        \item Il sistema aggiunge un nuovo marcatore di elenco (punto o numero successivo) alla riga seguente.
    \end{enumerate}
    \item \textbf{Postcondizioni}: Un nuovo elemento vuoto viene aggiunto all'elenco corrente.
    \item \textbf{Trigger}: L'utente desidera aggiungere una nuova voce all'elenco.
\end{itemize}

% UC 35 - Incremento livello elemento elenco
\hypertarget{uc35}{}
\subsubsection{UC 35: Incremento livello elemento elenco}

% use case diagram
\begin{figure}[H]
    \centering
    \includegraphics[width=0.5\textwidth]{incremento_livello_elenco.png}
    \caption{UC 35 – Incremento livello elemento elenco}
\end{figure}

\begin{itemize}
    \item \textbf{Attore Principale}: Utente
    \item \textbf{Precondizioni}: Il cursore si trova su un elemento di un elenco esistente.
    \item \textbf{Scenario Principale}:
    \begin{enumerate}
        \item L'utente attiva il comando di indentazione (incremento livello).
        \item Il sistema sposta l'elemento corrente al livello di annidamento successivo.
    \end{enumerate}
    \item \textbf{Postcondizioni}: L'elemento dell'elenco risulta annidato (sotto-elenco) rispetto al livello precedente.
    \item \textbf{Trigger}: L'utente desidera creare un sotto-elemento nell'elenco.
\end{itemize}

% UC 36 - Riduzione livello elemento elenco
\hypertarget{uc36}{}
\subsubsection{UC 36: Riduzione livello elemento elenco}

% use case diagram
\begin{figure}[H]
    \centering
    \includegraphics[width=0.5\textwidth]{riduzione_livello_elenco.png}
    \caption{UC 36 – Riduzione livello elemento elenco}
\end{figure}

\begin{itemize}
    \item \textbf{Attore Principale}: Utente
    \item \textbf{Precondizioni}: Il cursore si trova su un elemento annidato di un elenco.
    \item \textbf{Scenario Principale}:
    \begin{enumerate}
        \item L'utente attiva il comando di riduzione indentazione (riduzione livello).
        \item Il sistema riporta l'elemento corrente al livello di annidamento precedente.
    \end{enumerate}
    \item \textbf{Postcondizioni}: Il livello di annidamento dell'elemento risulta ridotto.
    \item \textbf{Trigger}: L'utente desidera riportare un sotto-elemento a un livello gerarchico superiore.
\end{itemize}

% UC 37 - Modifica elenco
\hypertarget{uc37}{}
\subsubsection{UC 37: Modifica elenco}

% use case diagram
\begin{figure}[H]
    \centering
    \includegraphics[width=0.5\textwidth]{modifica_elenco.png}
    \caption{UC 37 – Modifica elenco}
\end{figure}

\begin{itemize}
    \item \textbf{Attore Principale}: Utente
    \item \textbf{Precondizioni}: È presente almeno un elenco nel documento.
    \item \textbf{Scenario Principale}:
    \begin{enumerate}
        \item L'utente seleziona un elenco esistente.
        \item L'utente modifica la tipologia dell'elenco.
        \item Il sistema aggiorna l'elenco.
    \end{enumerate}
    \item \textbf{Postcondizioni}: L'elenco presenta le nuove proprietà impostate.
    \item \textbf{Trigger}: L'utente desidera modificare un elenco esistente.
\end{itemize}

% UC 38 - Disattivazione modalità elenco
\hypertarget{uc38}{}
\subsubsection{UC 38: Disattivazione modalità elenco}

% use case diagram
\begin{figure}[H]
    \centering
    \includegraphics[width=0.5\textwidth]{disattivazione_elenco.png}
    \caption{UC 38 – Disattivazione modalità elenco}
\end{figure}

\begin{itemize}
    \item \textbf{Attore Principale}: Utente
    \item \textbf{Precondizioni}: La modalità di inserimento elenco è attiva.
    \item \textbf{Scenario Principale}:
    \begin{enumerate}
        \item L'utente disattiva il comando "Elenco".
        \item Il sistema disabilita la modalità di inserimento elenco.
    \end{enumerate}
    \item \textbf{Postcondizioni}: Il testo inserito successivamente non verrà formattato come elenco.
    \item \textbf{Trigger}: L'utente desidera interrompere la modalità di inserimento elenco.
\end{itemize}

% UC 39 - Rimozione formattazione elenco da testo selezionato
\hypertarget{uc39}{}
\subsubsection{UC 39: Rimozione formattazione elenco da testo selezionato}

% use case diagram
\begin{figure}[H]
    \centering
    \includegraphics[width=0.5\textwidth]{rimozione_elenco.png}
    \caption{UC 39 – Rimozione formattazione elenco da testo selezionato}
\end{figure}

\begin{itemize}
    \item \textbf{Attore Principale}: Utente
    \item \textbf{Precondizioni}: Una porzione di elenco è selezionata.
    \item \textbf{Scenario Principale}:
    \begin{enumerate}
        \item L'utente disattiva il comando "Elenco".
        \item Il sistema rimuove la formattazione elenco dal testo selezionato.
    \end{enumerate}
    \item \textbf{Postcondizioni}: Il testo selezionato non presenta più la formattazione elenco.
    \item \textbf{Trigger}: L'utente desidera rimuovere la formattazione elenco da testo esistente.
\end{itemize}

%---------------GESTIONE TABELLE--------------------

% UC40 - Inserimento tabella
\hypertarget{uc40}{}
\subsubsection{UC 40: Inserimento tabella}

\begin{figure}[H]
    \centering
    \includegraphics[width=0.5\textwidth]{inserimento_tabella.png}
    \caption{UC 40 – Inserimento tabella}
\end{figure}

\begin{itemize}
    \item \textbf{Attore Principale}: Utente
    \item \textbf{Precondizioni}: Una nota è aperta nell'editor.
    \item \textbf{Scenario Principale}:
    \begin{enumerate}
        \item L'utente attiva la funzione di inserimento tabella.
        \item Il sistema richiede il numero di righe e colonne.
        \item L'utente inserisce i parametri desiderati.
        \item Il sistema genera la struttura Markdown della tabella.
        \item Il sistema inserisce la tabella nell'editor.
    \end{enumerate}
    \item \textbf{Postcondizioni}: Nell'editor è presente una nuova tabella.
    \item \textbf{Trigger}: L'utente desidera inserire una tabella nella nota.
    \item \textbf{Estensioni}:
    \begin{itemize}
        \item \hyperlink{uc40.1}{UC 40.1: Visualizzazione errore numero righe tabella} $\rightarrow$ Condizione: numero di righe non valido al passo 3.
        \item \hyperlink{uc40.2}{UC 40.2: Visualizzazione errore numero colonne tabella} $\rightarrow$ Condizione: numero di colonne non valido al passo 3.
    \end{itemize}
\end{itemize}

% UC40.1 - Visualizzazione errore numero righe tabella
\hypertarget{uc40.1}{}
\subsubsection{UC 40.1: Visualizzazione errore numero righe tabella}

\begin{figure}[H]
    \centering
    \includegraphics[width=0.6\textwidth]{errore_righe_tabella.png}
    \caption{UC 40.1 – Visualizzazione errore numero righe tabella}
\end{figure}

\begin{itemize}
    \item \textbf{Attore Principale}: Utente
    \item \textbf{Precondizioni}: È in corso la creazione di una tabella.
    \item \textbf{Scenario Principale}:
    \begin{enumerate}
        \item Il sistema rileva un valore non valido per il numero di righe.
        \item Il sistema mostra un messaggio di errore all'utente.
    \end{enumerate}
    \item \textbf{Postcondizioni}: La tabella non viene generata.
    \item \textbf{Trigger}: Inserimento di un parametro non valido per le righe.
\end{itemize}

% UC40.2 - Visualizzazione errore numero colonne tabella
\hypertarget{uc40.2}{}
\subsubsection{UC 40.2: Visualizzazione errore numero colonne tabella}

\begin{figure}[H]
    \centering
    \includegraphics[width=0.6\textwidth]{errore_colonne_tabella.png}
    \caption{UC 40.2 – Visualizzazione errore numero colonne tabella}
\end{figure}

\begin{itemize}
    \item \textbf{Attore Principale}: Utente
    \item \textbf{Precondizioni}: È in corso la creazione di una tabella.
    \item \textbf{Scenario Principale}:
    \begin{enumerate}
        \item Il sistema rileva un valore non valido per il numero di colonne.
        \item Il sistema mostra un messaggio di errore all'utente.
    \end{enumerate}
    \item \textbf{Postcondizioni}: La tabella non viene generata.
    \item \textbf{Trigger}: Inserimento di un parametro non valido per le colonne.
\end{itemize}

% UC41 - Modifica contenuto cella tabella
\hypertarget{uc41}{}
\subsubsection{UC 41: Modifica contenuto cella tabella}

\begin{figure}[H]
    \centering
    \includegraphics[width=0.5\textwidth]{modifica_cella_tabella.png}
    \caption{UC 41 – Modifica contenuto cella tabella}
\end{figure}

\begin{itemize}
    \item \textbf{Attore Principale}: Utente
    \item \textbf{Precondizioni}: Nell'editor è presente una tabella.
    \item \textbf{Scenario Principale}:
    \begin{enumerate}
        \item L'utente seleziona una cella della tabella.
        \item L'utente inserisce o modifica il contenuto della cella.
        \item Il sistema aggiorna la struttura della tabella.
    \end{enumerate}
    \item \textbf{Postcondizioni}: Il contenuto della cella risulta aggiornato.
    \item \textbf{Trigger}: L'utente desidera modificare una cella della tabella.
\end{itemize}

% UC42 - Inserimento riga tabella
\hypertarget{uc42}{}
\subsubsection{UC 42: Inserimento riga tabella}

\begin{figure}[H]
    \centering
    \includegraphics[width=0.5\textwidth]{inserimento_riga_tabella.png}
    \caption{UC 42 – Inserimento riga tabella}
\end{figure}

\begin{itemize}
    \item \textbf{Attore Principale}: Utente
    \item \textbf{Precondizioni}: Nell'editor è presente una tabella.
    \item \textbf{Scenario Principale}:
    \begin{enumerate}
        \item L'utente seleziona una riga della tabella.
        \item L'utente richiede l'inserimento di una nuova riga.
        \item Il sistema aggiunge una nuova riga vuota alla tabella.
    \end{enumerate}
    \item \textbf{Postcondizioni}: La tabella contiene una nuova riga.
    \item \textbf{Trigger}: L'utente desidera aggiungere una riga alla tabella.
\end{itemize}

% UC43 - Eliminazione riga tabella
\hypertarget{uc43}{}
\subsubsection{UC 43: Eliminazione riga tabella}

\begin{figure}[H]
    \centering
    \includegraphics[width=0.5\textwidth]{eliminazione_riga_tabella.png}
    \caption{UC 43 – Eliminazione riga tabella}
\end{figure}

\begin{itemize}
    \item \textbf{Attore Principale}: Utente
    \item \textbf{Precondizioni}: Nell'editor è presente una tabella con almeno una riga.
    \item \textbf{Scenario Principale}:
    \begin{enumerate}
        \item L'utente seleziona una riga della tabella.
        \item L'utente richiede la rimozione della riga.
        \item Il sistema elimina la riga selezionata.
    \end{enumerate}
    \item \textbf{Postcondizioni}: La riga selezionata viene rimossa dalla tabella.
    \item \textbf{Trigger}: L'utente desidera eliminare una riga della tabella.
\end{itemize}

% UC44 - Inserimento colonna tabella
\hypertarget{uc44}{}
\subsubsection{UC 44: Inserimento colonna tabella}

\begin{figure}[H]
    \centering
    \includegraphics[width=0.5\textwidth]{inserimento_colonna_tabella.png}
    \caption{UC 44 – Inserimento colonna tabella}
\end{figure}

\begin{itemize}
    \item \textbf{Attore Principale}: Utente
    \item \textbf{Precondizioni}: Nell'editor è presente una tabella.
    \item \textbf{Scenario Principale}:
    \begin{enumerate}
        \item L'utente seleziona una colonna della tabella.
        \item L'utente richiede l'inserimento di una nuova colonna.
        \item Il sistema aggiunge una nuova colonna vuota alla tabella.
    \end{enumerate}
    \item \textbf{Postcondizioni}: La tabella contiene una nuova colonna.
    \item \textbf{Trigger}: L'utente desidera aggiungere una colonna alla tabella.
\end{itemize}

% UC45 - Eliminazione colonna tabella
\hypertarget{uc45}{}
\subsubsection{UC 45: Eliminazione colonna tabella}

\begin{figure}[H]
    \centering
    \includegraphics[width=0.5\textwidth]{eliminazione_colonna_tabella.png}
    \caption{UC 45 – Eliminazione colonna tabella}
\end{figure}

\begin{itemize}
    \item \textbf{Attore Principale}: Utente
    \item \textbf{Precondizioni}: Nell'editor è presente una tabella con almeno una colonna.
    \item \textbf{Scenario Principale}:
    \begin{enumerate}
        \item L'utente seleziona una colonna della tabella.
        \item L'utente richiede la rimozione della colonna.
        \item Il sistema elimina la colonna selezionata.
    \end{enumerate}
    \item \textbf{Postcondizioni}: La colonna selezionata viene rimossa dalla tabella.
    \item \textbf{Trigger}: L'utente desidera eliminare una colonna della tabella.
\end{itemize}

% UC46 - Eliminazione tabella
\hypertarget{uc46}{}
\subsubsection{UC 46: Eliminazione tabella}

\begin{figure}[H]
    \centering
    \includegraphics[width=0.5\textwidth]{eliminazione_tabella.png}
    \caption{UC 46 – Eliminazione tabella}
\end{figure}

\begin{itemize}
    \item \textbf{Attore Principale}: Utente
    \item \textbf{Precondizioni}: Nell'editor è presente una tabella.
    \item \textbf{Scenario Principale}:
    \begin{enumerate}
        \item L'utente seleziona la tabella.
        \item L'utente richiede la rimozione della tabella.
        \item Il sistema elimina la tabella dall'editor.
    \end{enumerate}
    \item \textbf{Postcondizioni}: La tabella viene rimossa dalla nota.
    \item \textbf{Trigger}: L'utente desidera eliminare una tabella.
\end{itemize}

%---------------GESTIONE MODIFICHE------------------

% UC47 - Azione di Undo
\hypertarget{uc47}{}
\subsubsection{UC 47: Azione di Undo\textsubscript{G}}

% use case diagram
\begin{figure}[H]
    \centering
    \includegraphics[width=0.5\textwidth]{azione_undo.png}
    \caption{UC 47 – Azione di Undo}
\end{figure}

\begin{itemize}
    \item \textbf{Attore Principale}: Utente
    \item \textbf{Precondizioni}: Esiste almeno una modifica annullabile.
    \item \textbf{Scenario Principale}:
    \begin{enumerate}
        \item L'utente attiva il comando "Undo".
        \item Il sistema annulla l'ultima modifica effettuata.
    \end{enumerate}
    \item \textbf{Postcondizioni}: Il documento viene ripristinato allo stato precedente.
    \item \textbf{Trigger}: L'utente desidera annullare un'operazione eseguita.
\end{itemize}

% UC48 - Azione di Redo
\hypertarget{uc48}{}
\subsubsection{UC 48: Azione di Redo\textsubscript{G}}

% use case diagram
\begin{figure}[H]
    \centering
    \includegraphics[width=0.5\textwidth]{azione_redo.png}
    \caption{UC 48 – Azione di Redo}
\end{figure}

\begin{itemize}
    \item \textbf{Attore Principale}: Utente
    \item \textbf{Precondizioni}: È stata precedentemente eseguita almeno un'operazione di Undo.
    \item \textbf{Scenario Principale}:
    \begin{enumerate}
        \item L'utente attiva il comando "Redo".
        \item Il sistema ripristina l'ultima modifica annullata.
    \end{enumerate}
    \item \textbf{Postcondizioni}: La modifica precedentemente annullata viene riapplicata.
    \item \textbf{Trigger}: L'utente desidera ripristinare una modifica annullata.
\end{itemize}

%---------------GESTIONE VISUALIZZAZIONE----------------

% UC49 - Visualizzazione esclusiva Editor
\hypertarget{uc49}{}
\subsubsection{UC 49: Visualizzazione esclusiva Editor}

% use case diagram
\begin{figure}[H]
    \centering
    \includegraphics[width=0.5\textwidth]{visualizzazione_editor.png}
    \caption{UC 49 – Visualizzazione esclusiva Editor}
\end{figure}

\begin{itemize}
    \item \textbf{Attore Principale}: Utente
    \item \textbf{Precondizioni}: L'interfaccia dell'applicazione è attiva.
    \item \textbf{Scenario Principale}:
    \begin{enumerate}
        \item L'utente attiva la modalità "Solo Editor\textsubscript{G}".
        \item Il sistema mostra esclusivamente l'area di modifica del documento.
    \end{enumerate}
    \item \textbf{Postcondizioni}: L'interfaccia mostra unicamente l'editor testuale.
    \item \textbf{Trigger}: L'utente desidera concentrarsi esclusivamente sulla scrittura del documento.
\end{itemize}

% UC 50 - Visualizzazione esclusiva Render\textsubscript{G}
\hypertarget{uc50}{}
\subsubsection{UC 50: Visualizzazione esclusiva Render}

% use case diagram
\begin{figure}[H]
    \centering
    \includegraphics[width=0.5\textwidth]{visualizzazione_render.png}
    \caption{UC 50 – Visualizzazione esclusiva Render}
\end{figure}

\begin{itemize}
    \item \textbf{Attore Principale}: Utente
    \item \textbf{Precondizioni}: L'interfaccia dell'applicazione è attiva.
    \item \textbf{Scenario Principale}:
    \begin{enumerate}
        \item L'utente attiva la modalità "Solo Render\textsubscript{G}".
        \item Il sistema mostra esclusivamente l'anteprima formattata del documento.
    \end{enumerate}
    \item \textbf{Postcondizioni}: L'interfaccia mostra unicamente l'anteprima del documento renderizzato.
    \item \textbf{Trigger}: L'utente desidera visualizzare il risultato finale del documento.
\end{itemize}

% UC 51 - Visualizzazione Split View\textsubscript{G}
\hypertarget{uc51}{}
\subsubsection{UC 51: Visualizzazione Split View}

% use case diagram
\begin{figure}[H]
    \centering
    \includegraphics[width=0.5\textwidth]{visualizzazione_split_view.png}
    \caption{UC 51 – Visualizzazione Split View}
\end{figure}

\begin{itemize}
    \item \textbf{Attore Principale}: Utente
    \item \textbf{Precondizioni}: L'interfaccia dell'applicazione è attiva.
    \item \textbf{Scenario Principale}:
    \begin{enumerate}
        \item L'utente attiva la modalità "Split View".
        \item Il sistema mostra contemporaneamente l'editor e l'anteprima del documento.
    \end{enumerate}
    \item \textbf{Postcondizioni}: L'interfaccia mostra simultaneamente l'editor testuale e l'anteprima renderizzata.
    \item \textbf{Trigger}: L'utente desidera visualizzare il rendering del documento durante la modifica del contenuto.
\end{itemize}

% --------------------------------------------------------------------------
% 2. Macroarea 2: Interazione AI (LLM)
% --------------------------------------------------------------------------
\newpage
\subsection{Macroarea 2 — Interazione AI (LLM)}
\textit{Elaborazione del testo tramite modelli linguistici e gestione degli output generati.}

\noindent
Per tutte le funzionalità basate su LLM, il sistema visualizza durante l'elaborazione
un indicatore di caricamento e rende disponibile un comando di interruzione.
Al completamento, la proposta generata viene mostrata in una finestra modale dedicata,
dalla quale l'utente può accettarla, rifiutarla o richiederne la rigenerazione.
Nessuna proposta viene applicata automaticamente al contenuto della nota.


% UC52 - Richiesta riassunto
\hypertarget{uc52}{}
\subsubsection{UC 52: Richiesta riassunto}

\begin{figure}[H]
    \centering
    \includegraphics[width=0.6\textwidth]{richiesta_riassunto.png}
    \caption{UC 52 – Richiesta riassunto}
\end{figure}

\begin{itemize}
    \item \textbf{Attore Principale}: Utente
    \item \textbf{Attore Secondario}: LLM
    \item \textbf{Precondizioni}: Una porzione di testo è selezionata nell'editor e il servizio AI è disponibile.
    \item \textbf{Scenario Principale}:
    \begin{enumerate}
        \item L'utente richiede la generazione di un riassunto del testo selezionato.
        \item Il sistema acquisisce la richiesta dell'utente.
    \end{enumerate}
    \item \textbf{Postcondizioni}: La richiesta di riassunto è stata acquisita dal sistema.
    \item \textbf{Trigger}: L'utente desidera ottenere una sintesi del testo selezionato.
\end{itemize}


% UC53 - Generazione riassunto
\hypertarget{uc53}{}
\subsubsection{UC 53: Generazione riassunto}

\begin{figure}[H]
    \centering
    \includegraphics[width=0.6\textwidth]{generazione_riassunto.png}
    \caption{UC 53 – Generazione riassunto}
\end{figure}

\begin{itemize}
    \item \textbf{Attore Principale}: Utente
    \item \textbf{Attore Secondario}: LLM
    \item \textbf{Precondizioni}: La richiesta di riassunto è stata acquisita dal sistema.
    \item \textbf{Scenario Principale}:
    \begin{enumerate}
        \item Il sistema inoltra al servizio LLM la richiesta di riassunto.
        \item Il servizio LLM prende in carico la richiesta e avvia l'elaborazione.
    \end{enumerate}
    \item \textbf{Postcondizioni}: La richiesta di riassunto è in elaborazione presso il servizio LLM.
    \item \textbf{Trigger}: Il sistema ha acquisito la richiesta di riassunto.
\end{itemize}


% UC54 - Visualizzazione elaborazione riassunto
\hypertarget{uc54}{}
\subsubsection{UC 54: Visualizzazione elaborazione riassunto}

\begin{figure}[H]
    \centering
    \includegraphics[width=0.6\textwidth]{elaborazione_riassunto.png}
    \caption{UC 54 – Visualizzazione elaborazione riassunto}
\end{figure}

\begin{itemize}
    \item \textbf{Attore Principale}: Utente
    \item \textbf{Precondizioni}: La richiesta di riassunto è in elaborazione presso il servizio LLM.
    \item \textbf{Scenario Principale}:
    \begin{enumerate}
        \item Il sistema visualizza una finestra modale contenente un indicatore di caricamento.
        \item Il sistema rende disponibile il comando ``Interrompi''.
    \end{enumerate}
    \item \textbf{Postcondizioni}: La finestra modale di elaborazione è visualizzata e la richiesta è ancora in corso.
    \item \textbf{Trigger}: Il servizio LLM ha avviato l'elaborazione della richiesta di riassunto.
\end{itemize}


% UC55 - Interruzione elaborazione riassunto
\hypertarget{uc55}{}
\subsubsection{UC 55: Interruzione elaborazione riassunto}

\begin{figure}[H]
    \centering
    \includegraphics[width=0.6\textwidth]{interruzione_riassunto.png}
    \caption{UC 55 – Interruzione elaborazione riassunto}
\end{figure}

\begin{itemize}
    \item \textbf{Attore Principale}: Utente
    \item \textbf{Attore Secondario}: LLM
    \item \textbf{Precondizioni}: La finestra modale di elaborazione è visualizzata e la richiesta è ancora in corso.
    \item \textbf{Scenario Principale}:
    \begin{enumerate}
        \item L'utente attiva il comando ``Interrompi''.
        \item Il sistema richiede al servizio LLM l'interruzione dell'elaborazione.
        \item Il sistema chiude la finestra modale di elaborazione.
    \end{enumerate}
    \item \textbf{Postcondizioni}: L'elaborazione è interrotta e il contenuto della nota rimane invariato.
    \item \textbf{Trigger}: L'utente desidera interrompere la richiesta ancora in elaborazione.
\end{itemize}


% UC56 - Visualizzazione proposta riassunto
\hypertarget{uc56}{}
\subsubsection{UC 56: Visualizzazione proposta riassunto}

\begin{figure}[H]
    \centering
    \includegraphics[width=0.6\textwidth]{proposta_riassunto.png}
    \caption{UC 56 – Visualizzazione proposta riassunto}
\end{figure}

\begin{itemize}
    \item \textbf{Attore Principale}: Utente
    \item \textbf{Precondizioni}: Il servizio LLM ha completato la generazione della proposta di riassunto.
    \item \textbf{Scenario Principale}:
    \begin{enumerate}
        \item Il sistema chiude la finestra modale di elaborazione.
        \item Il sistema visualizza il riassunto del testo in una finestra modale dedicata.
        \item Il sistema rende disponibili i comandi ``Accetta'', ``Rifiuta'' e ``Rigenera''.
    \end{enumerate}
    \item \textbf{Postcondizioni}: La proposta è visualizzata e il contenuto della nota rimane invariato.
    \item \textbf{Trigger}: Il servizio LLM ha restituito la proposta di riassunto.
\end{itemize}


% UC57 - Accettazione proposta riassunto
\hypertarget{uc57}{}
\subsubsection{UC 57: Accettazione proposta riassunto}

\begin{figure}[H]
    \centering
    \includegraphics[width=0.6\textwidth]{accettazione_riassunto.png}
    \caption{UC 57 – Accettazione proposta riassunto}
\end{figure}

\begin{itemize}
    \item \textbf{Attore Principale}: Utente
    \item \textbf{Precondizioni}: La finestra modale contenente la proposta di riassunto è visualizzata.
    \item \textbf{Scenario Principale}:
    \begin{enumerate}
        \item L'utente attiva il comando ``Accetta''.
        \item Il sistema sostituisce il testo inizialmente selezionato con il riassunto del testo proposto.
        \item Il sistema chiude la finestra modale della proposta.
    \end{enumerate}
    \item \textbf{Postcondizioni}: Il testo inizialmente selezionato è sostituito con il riassunto del testo accettato e la finestra modale è chiusa.
    \item \textbf{Trigger}: L'utente desidera applicare la proposta visualizzata.
\end{itemize}


% UC58 - Rifiuto proposta riassunto
\hypertarget{uc58}{}
\subsubsection{UC 58: Rifiuto proposta riassunto}

\begin{figure}[H]
    \centering
    \includegraphics[width=0.6\textwidth]{rifiuto_riassunto.png}
    \caption{UC 58 – Rifiuto proposta riassunto}
\end{figure}

\begin{itemize}
    \item \textbf{Attore Principale}: Utente
    \item \textbf{Precondizioni}: La finestra modale contenente la proposta di riassunto è visualizzata.
    \item \textbf{Scenario Principale}:
    \begin{enumerate}
        \item L'utente attiva il comando ``Rifiuta''.
        \item Il sistema scarta la proposta visualizzata.
        \item Il sistema chiude la finestra modale della proposta.
    \end{enumerate}
    \item \textbf{Postcondizioni}: La proposta è scartata, la finestra modale è chiusa e il contenuto della nota rimane invariato.
    \item \textbf{Trigger}: L'utente non desidera applicare la proposta visualizzata.
\end{itemize}


% UC59 - Rigenerazione proposta riassunto
\hypertarget{uc59}{}
\subsubsection{UC 59: Rigenerazione proposta riassunto}

\begin{figure}[H]
    \centering
    \includegraphics[width=0.6\textwidth]{rigenerazione_riassunto.png}
    \caption{UC 59 – Rigenerazione proposta riassunto}
\end{figure}

\begin{itemize}
    \item \textbf{Attore Principale}: Utente
    \item \textbf{Attore Secondario}: LLM
    \item \textbf{Precondizioni}: La finestra modale contenente la proposta di riassunto è visualizzata.
    \item \textbf{Scenario Principale}:
    \begin{enumerate}
        \item L'utente attiva il comando ``Rigenera''.
        \item Il sistema scarta la proposta visualizzata.
        \item Il sistema inoltra al servizio LLM una nuova richiesta di riassunto.
        \item Il sistema visualizza nuovamente la finestra modale di elaborazione.
    \end{enumerate}
    \item \textbf{Postcondizioni}: Una nuova richiesta è in elaborazione e il contenuto della nota rimane invariato.
    \item \textbf{Trigger}: L'utente desidera ottenere una proposta differente.
\end{itemize}


% UC60 - Visualizzazione errore generazione riassunto
\hypertarget{uc60}{}
\subsubsection{UC 60: Visualizzazione errore generazione riassunto}

\begin{figure}[H]
    \centering
    \includegraphics[width=0.6\textwidth]{errore_riassunto.png}
    \caption{UC 60 – Visualizzazione errore generazione riassunto}
\end{figure}

\begin{itemize}
    \item \textbf{Attore Principale}: Utente
    \item \textbf{Precondizioni}: È in corso una richiesta di riassunto.
    \item \textbf{Scenario Principale}:
    \begin{enumerate}
        \item Si verifica un errore durante la comunicazione con il servizio LLM o durante l'elaborazione della richiesta.
        \item Il sistema chiude la finestra modale di elaborazione.
        \item Il sistema visualizza un messaggio che segnala il mancato completamento della richiesta.
    \end{enumerate}
    \item \textbf{Postcondizioni}: L'utente è informato dell'errore e il contenuto della nota rimane invariato.
    \item \textbf{Trigger}: Si verifica un errore prima del completamento della richiesta di riassunto.
\end{itemize}

% __________________________________________________%


% UC61 - Richiesta traduzione
\hypertarget{uc61}{}
\subsubsection{UC 61: Richiesta traduzione}

\begin{figure}[H]
    \centering
    \includegraphics[width=0.6\textwidth]{richiesta_traduzione.png}
    \caption{UC 61 – Richiesta traduzione}
\end{figure}

\begin{itemize}
    \item \textbf{Attore Principale}: Utente
    \item \textbf{Attore Secondario}: LLM
    \item \textbf{Precondizioni}: Una porzione di testo è selezionata nell'editor e il servizio AI è disponibile.
    \item \textbf{Scenario Principale}:
    \begin{enumerate}
        \item L'utente richiede la generazione di una traduzione del testo selezionato.
        \item Il sistema acquisisce la richiesta dell'utente.
    \end{enumerate}
    \item \textbf{Postcondizioni}: La richiesta di traduzione è stata acquisita dal sistema.
    \item \textbf{Trigger}: L'utente desidera ottenere una traduzione del testo selezionato.
\end{itemize}


% UC62 - Generazione traduzione
\hypertarget{uc62}{}
\subsubsection{UC 62: Generazione traduzione}

\begin{figure}[H]
    \centering
    \includegraphics[width=0.6\textwidth]{generazione_traduzione.png}
    \caption{UC 62 – Generazione traduzione}
\end{figure}

\begin{itemize}
    \item \textbf{Attore Principale}: Utente
    \item \textbf{Attore Secondario}: LLM
    \item \textbf{Precondizioni}: La richiesta di traduzione è stata acquisita dal sistema.
    \item \textbf{Scenario Principale}:
    \begin{enumerate}
        \item Il sistema inoltra al servizio LLM la richiesta di traduzione.
        \item Il servizio LLM prende in carico la richiesta e avvia l'elaborazione.
    \end{enumerate}
    \item \textbf{Postcondizioni}: La richiesta di traduzione è in elaborazione presso il servizio LLM.
    \item \textbf{Trigger}: Il sistema ha acquisito la richiesta di traduzione.
\end{itemize}


% UC63 - Visualizzazione elaborazione traduzione
\hypertarget{uc63}{}
\subsubsection{UC 63: Visualizzazione elaborazione traduzione}

\begin{figure}[H]
    \centering
    \includegraphics[width=0.6\textwidth]{elaborazione_traduzione.png}
    \caption{UC 63 – Visualizzazione elaborazione traduzione}
\end{figure}

\begin{itemize}
    \item \textbf{Attore Principale}: Utente
    \item \textbf{Precondizioni}: La richiesta di traduzione è in elaborazione presso il servizio LLM.
    \item \textbf{Scenario Principale}:
    \begin{enumerate}
        \item Il sistema visualizza una finestra modale contenente un indicatore di caricamento.
        \item Il sistema rende disponibile il comando ``Interrompi''.
    \end{enumerate}
    \item \textbf{Postcondizioni}: La finestra modale di elaborazione è visualizzata e la richiesta è ancora in corso.
    \item \textbf{Trigger}: Il servizio LLM ha avviato l'elaborazione della richiesta di traduzione.
\end{itemize}


% UC64 - Interruzione elaborazione traduzione
\hypertarget{uc64}{}
\subsubsection{UC 64: Interruzione elaborazione traduzione}

\begin{figure}[H]
    \centering
    \includegraphics[width=0.6\textwidth]{interruzione_traduzione.png}
    \caption{UC 64 – Interruzione elaborazione traduzione}
\end{figure}

\begin{itemize}
    \item \textbf{Attore Principale}: Utente
    \item \textbf{Attore Secondario}: LLM
    \item \textbf{Precondizioni}: La finestra modale di elaborazione è visualizzata e la richiesta è ancora in corso.
    \item \textbf{Scenario Principale}:
    \begin{enumerate}
        \item L'utente attiva il comando ``Interrompi''.
        \item Il sistema richiede al servizio LLM l'interruzione dell'elaborazione.
        \item Il sistema chiude la finestra modale di elaborazione.
    \end{enumerate}
    \item \textbf{Postcondizioni}: L'elaborazione è interrotta e il contenuto della nota rimane invariato.
    \item \textbf{Trigger}: L'utente desidera interrompere la richiesta ancora in elaborazione.
\end{itemize}


% UC65 - Visualizzazione proposta traduzione
\hypertarget{uc65}{}
\subsubsection{UC 65: Visualizzazione proposta traduzione}

\begin{figure}[H]
    \centering
    \includegraphics[width=0.6\textwidth]{proposta_traduzione.png}
    \caption{UC 65 – Visualizzazione proposta traduzione}
\end{figure}

\begin{itemize}
    \item \textbf{Attore Principale}: Utente
    \item \textbf{Precondizioni}: Il servizio LLM ha completato la generazione della proposta di traduzione.
    \item \textbf{Scenario Principale}:
    \begin{enumerate}
        \item Il sistema chiude la finestra modale di elaborazione.
        \item Il sistema visualizza la traduzione del testo in una finestra modale dedicata.
        \item Il sistema rende disponibili i comandi ``Accetta'', ``Rifiuta'' e ``Rigenera''.
    \end{enumerate}
    \item \textbf{Postcondizioni}: La proposta è visualizzata e il contenuto della nota rimane invariato.
    \item \textbf{Trigger}: Il servizio LLM ha restituito la proposta di traduzione.
\end{itemize}


% UC66 - Accettazione proposta traduzione
\hypertarget{uc66}{}
\subsubsection{UC 66: Accettazione proposta traduzione}

\begin{figure}[H]
    \centering
    \includegraphics[width=0.6\textwidth]{accettazione_traduzione.png}
    \caption{UC 66 – Accettazione proposta traduzione}
\end{figure}

\begin{itemize}
    \item \textbf{Attore Principale}: Utente
    \item \textbf{Precondizioni}: La finestra modale contenente la proposta di traduzione è visualizzata.
    \item \textbf{Scenario Principale}:
    \begin{enumerate}
        \item L'utente attiva il comando ``Accetta''.
        \item Il sistema sostituisce il testo inizialmente selezionato con la traduzione del testo proposto.
        \item Il sistema chiude la finestra modale della proposta.
    \end{enumerate}
    \item \textbf{Postcondizioni}: Il testo inizialmente selezionato è sostituito con la traduzione del testo accettata e la finestra modale è chiusa.
    \item \textbf{Trigger}: L'utente desidera applicare la proposta visualizzata.
\end{itemize}


% UC67 - Rifiuto proposta traduzione
\hypertarget{uc67}{}
\subsubsection{UC 67: Rifiuto proposta traduzione}

\begin{figure}[H]
    \centering
    \includegraphics[width=0.6\textwidth]{rifiuto_traduzione.png}
    \caption{UC 67 – Rifiuto proposta traduzione}
\end{figure}

\begin{itemize}
    \item \textbf{Attore Principale}: Utente
    \item \textbf{Precondizioni}: La finestra modale contenente la proposta di traduzione è visualizzata.
    \item \textbf{Scenario Principale}:
    \begin{enumerate}
        \item L'utente attiva il comando ``Rifiuta''.
        \item Il sistema scarta la proposta visualizzata.
        \item Il sistema chiude la finestra modale della proposta.
    \end{enumerate}
    \item \textbf{Postcondizioni}: La proposta è scartata, la finestra modale è chiusa e il contenuto della nota rimane invariato.
    \item \textbf{Trigger}: L'utente non desidera applicare la proposta visualizzata.
\end{itemize}


% UC68 - Rigenerazione proposta traduzione
\hypertarget{uc68}{}
\subsubsection{UC 68: Rigenerazione proposta traduzione}

\begin{figure}[H]
    \centering
    \includegraphics[width=0.6\textwidth]{rigenerazione_traduzione.png}
    \caption{UC 68 – Rigenerazione proposta traduzione}
\end{figure}

\begin{itemize}
    \item \textbf{Attore Principale}: Utente
    \item \textbf{Attore Secondario}: LLM
    \item \textbf{Precondizioni}: La finestra modale contenente la proposta di traduzione è visualizzata.
    \item \textbf{Scenario Principale}:
    \begin{enumerate}
        \item L'utente attiva il comando ``Rigenera''.
        \item Il sistema scarta la proposta visualizzata.
        \item Il sistema inoltra al servizio LLM una nuova richiesta di traduzione.
        \item Il sistema visualizza nuovamente la finestra modale di elaborazione.
    \end{enumerate}
    \item \textbf{Postcondizioni}: Una nuova richiesta è in elaborazione e il contenuto della nota rimane invariato.
    \item \textbf{Trigger}: L'utente desidera ottenere una proposta differente.
\end{itemize}


% UC69 - Visualizzazione errore generazione traduzione
\hypertarget{uc69}{}
\subsubsection{UC 69: Visualizzazione errore generazione traduzione}

\begin{figure}[H]
    \centering
    \includegraphics[width=0.6\textwidth]{errore_traduzione.png}
    \caption{UC 69 – Visualizzazione errore generazione traduzione}
\end{figure}

\begin{itemize}
    \item \textbf{Attore Principale}: Utente
    \item \textbf{Precondizioni}: È in corso una richiesta di traduzione.
    \item \textbf{Scenario Principale}:
    \begin{enumerate}
        \item Si verifica un errore durante la comunicazione con il servizio LLM o durante l'elaborazione della richiesta.
        \item Il sistema chiude la finestra modale di elaborazione.
        \item Il sistema visualizza un messaggio che segnala il mancato completamento della richiesta.
    \end{enumerate}
    \item \textbf{Postcondizioni}: L'utente è informato dell'errore e il contenuto della nota rimane invariato.
    \item \textbf{Trigger}: Si verifica un errore prima del completamento della richiesta di traduzione.
\end{itemize}

% __________________________________________________%


% UC70 - Richiesta riscrittura
\hypertarget{uc70}{}
\subsubsection{UC 70: Richiesta riscrittura}

\begin{figure}[H]
    \centering
    \includegraphics[width=0.6\textwidth]{richiesta_riscrittura.png}
    \caption{UC 70 – Richiesta riscrittura}
\end{figure}

\begin{itemize}
    \item \textbf{Attore Principale}: Utente
    \item \textbf{Attore Secondario}: LLM
    \item \textbf{Precondizioni}: Una porzione di testo è selezionata nell'editor e il servizio AI è disponibile.
    \item \textbf{Scenario Principale}:
    \begin{enumerate}
        \item L'utente richiede la generazione di una riscrittura del testo selezionato.
        \item Il sistema acquisisce la richiesta dell'utente.
    \end{enumerate}
    \item \textbf{Postcondizioni}: La richiesta di riscrittura è stata acquisita dal sistema.
    \item \textbf{Trigger}: L'utente desidera ottenere una riscrittura del testo selezionato.
\end{itemize}


% UC71 - Generazione riscrittura
\hypertarget{uc71}{}
\subsubsection{UC 71: Generazione riscrittura}

\begin{figure}[H]
    \centering
    \includegraphics[width=0.6\textwidth]{generazione_riscrittura.png}
    \caption{UC 71 – Generazione riscrittura}
\end{figure}

\begin{itemize}
    \item \textbf{Attore Principale}: Utente
    \item \textbf{Attore Secondario}: LLM
    \item \textbf{Precondizioni}: La richiesta di riscrittura è stata acquisita dal sistema.
    \item \textbf{Scenario Principale}:
    \begin{enumerate}
        \item Il sistema inoltra al servizio LLM la richiesta di riscrittura.
        \item Il servizio LLM prende in carico la richiesta e avvia l'elaborazione.
    \end{enumerate}
    \item \textbf{Postcondizioni}: La richiesta di riscrittura è in elaborazione presso il servizio LLM.
    \item \textbf{Trigger}: Il sistema ha acquisito la richiesta di riscrittura.
\end{itemize}


% UC72 - Visualizzazione elaborazione riscrittura
\hypertarget{uc72}{}
\subsubsection{UC 72: Visualizzazione elaborazione riscrittura}

\begin{figure}[H]
    \centering
    \includegraphics[width=0.6\textwidth]{elaborazione_riscrittura.png}
    \caption{UC 72 – Visualizzazione elaborazione riscrittura}
\end{figure}

\begin{itemize}
    \item \textbf{Attore Principale}: Utente
    \item \textbf{Precondizioni}: La richiesta di riscrittura è in elaborazione presso il servizio LLM.
    \item \textbf{Scenario Principale}:
    \begin{enumerate}
        \item Il sistema visualizza una finestra modale contenente un indicatore di caricamento.
        \item Il sistema rende disponibile il comando ``Interrompi''.
    \end{enumerate}
    \item \textbf{Postcondizioni}: La finestra modale di elaborazione è visualizzata e la richiesta è ancora in corso.
    \item \textbf{Trigger}: Il servizio LLM ha avviato l'elaborazione della richiesta di riscrittura.
\end{itemize}


% UC73 - Interruzione elaborazione riscrittura
\hypertarget{uc73}{}
\subsubsection{UC 73: Interruzione elaborazione riscrittura}

\begin{figure}[H]
    \centering
    \includegraphics[width=0.6\textwidth]{interruzione_riscrittura.png}
    \caption{UC 73 – Interruzione elaborazione riscrittura}
\end{figure}

\begin{itemize}
    \item \textbf{Attore Principale}: Utente
    \item \textbf{Attore Secondario}: LLM
    \item \textbf{Precondizioni}: La finestra modale di elaborazione è visualizzata e la richiesta è ancora in corso.
    \item \textbf{Scenario Principale}:
    \begin{enumerate}
        \item L'utente attiva il comando ``Interrompi''.
        \item Il sistema richiede al servizio LLM l'interruzione dell'elaborazione.
        \item Il sistema chiude la finestra modale di elaborazione.
    \end{enumerate}
    \item \textbf{Postcondizioni}: L'elaborazione è interrotta e il contenuto della nota rimane invariato.
    \item \textbf{Trigger}: L'utente desidera interrompere la richiesta ancora in elaborazione.
\end{itemize}


% UC74 - Visualizzazione proposta riscrittura
\hypertarget{uc74}{}
\subsubsection{UC 74: Visualizzazione proposta riscrittura}

\begin{figure}[H]
    \centering
    \includegraphics[width=0.6\textwidth]{proposta_riscrittura.png}
    \caption{UC 74 – Visualizzazione proposta riscrittura}
\end{figure}

\begin{itemize}
    \item \textbf{Attore Principale}: Utente
    \item \textbf{Precondizioni}: Il servizio LLM ha completato la generazione della proposta di riscrittura.
    \item \textbf{Scenario Principale}:
    \begin{enumerate}
        \item Il sistema chiude la finestra modale di elaborazione.
        \item Il sistema visualizza la riscrittura del testo in una finestra modale dedicata.
        \item Il sistema rende disponibili i comandi ``Accetta'', ``Rifiuta'' e ``Rigenera''.
    \end{enumerate}
    \item \textbf{Postcondizioni}: La proposta è visualizzata e il contenuto della nota rimane invariato.
    \item \textbf{Trigger}: Il servizio LLM ha restituito la proposta di riscrittura.
\end{itemize}


% UC75 - Accettazione proposta riscrittura
\hypertarget{uc75}{}
\subsubsection{UC 75: Accettazione proposta riscrittura}

\begin{figure}[H]
    \centering
    \includegraphics[width=0.6\textwidth]{accettazione_riscrittura.png}
    \caption{UC 75 – Accettazione proposta riscrittura}
\end{figure}

\begin{itemize}
    \item \textbf{Attore Principale}: Utente
    \item \textbf{Precondizioni}: La finestra modale contenente la proposta di riscrittura è visualizzata.
    \item \textbf{Scenario Principale}:
    \begin{enumerate}
        \item L'utente attiva il comando ``Accetta''.
        \item Il sistema sostituisce il testo inizialmente selezionato con la riscrittura del testo proposto.
        \item Il sistema chiude la finestra modale della proposta.
    \end{enumerate}
    \item \textbf{Postcondizioni}: Il testo inizialmente selezionato è sostituito con la riscrittura del testo accettata e la finestra modale è chiusa.
    \item \textbf{Trigger}: L'utente desidera applicare la proposta visualizzata.
\end{itemize}


% UC76 - Rifiuto proposta riscrittura
\hypertarget{uc76}{}
\subsubsection{UC 76: Rifiuto proposta riscrittura}

\begin{figure}[H]
    \centering
    \includegraphics[width=0.6\textwidth]{rifiuto_riscrittura.png}
    \caption{UC 76 – Rifiuto proposta riscrittura}
\end{figure}

\begin{itemize}
    \item \textbf{Attore Principale}: Utente
    \item \textbf{Precondizioni}: La finestra modale contenente la proposta di riscrittura è visualizzata.
    \item \textbf{Scenario Principale}:
    \begin{enumerate}
        \item L'utente attiva il comando ``Rifiuta''.
        \item Il sistema scarta la proposta visualizzata.
        \item Il sistema chiude la finestra modale della proposta.
    \end{enumerate}
    \item \textbf{Postcondizioni}: La proposta è scartata, la finestra modale è chiusa e il contenuto della nota rimane invariato.
    \item \textbf{Trigger}: L'utente non desidera applicare la proposta visualizzata.
\end{itemize}


% UC77 - Rigenerazione proposta riscrittura
\hypertarget{uc77}{}
\subsubsection{UC 77: Rigenerazione proposta riscrittura}

\begin{figure}[H]
    \centering
    \includegraphics[width=0.6\textwidth]{rigenerazione_riscrittura.png}
    \caption{UC 77 – Rigenerazione proposta riscrittura}
\end{figure}

\begin{itemize}
    \item \textbf{Attore Principale}: Utente
    \item \textbf{Attore Secondario}: LLM
    \item \textbf{Precondizioni}: La finestra modale contenente la proposta di riscrittura è visualizzata.
    \item \textbf{Scenario Principale}:
    \begin{enumerate}
        \item L'utente attiva il comando ``Rigenera''.
        \item Il sistema scarta la proposta visualizzata.
        \item Il sistema inoltra al servizio LLM una nuova richiesta di riscrittura.
        \item Il sistema visualizza nuovamente la finestra modale di elaborazione.
    \end{enumerate}
    \item \textbf{Postcondizioni}: Una nuova richiesta è in elaborazione e il contenuto della nota rimane invariato.
    \item \textbf{Trigger}: L'utente desidera ottenere una proposta differente.
\end{itemize}


% UC78 - Visualizzazione errore generazione riscrittura
\hypertarget{uc78}{}
\subsubsection{UC 78: Visualizzazione errore generazione riscrittura}

\begin{figure}[H]
    \centering
    \includegraphics[width=0.6\textwidth]{errore_riscrittura.png}
    \caption{UC 78 – Visualizzazione errore generazione riscrittura}
\end{figure}

\begin{itemize}
    \item \textbf{Attore Principale}: Utente
    \item \textbf{Precondizioni}: È in corso una richiesta di riscrittura.
    \item \textbf{Scenario Principale}:
    \begin{enumerate}
        \item Si verifica un errore durante la comunicazione con il servizio LLM o durante l'elaborazione della richiesta.
        \item Il sistema chiude la finestra modale di elaborazione.
        \item Il sistema visualizza un messaggio che segnala il mancato completamento della richiesta.
    \end{enumerate}
    \item \textbf{Postcondizioni}: L'utente è informato dell'errore e il contenuto della nota rimane invariato.
    \item \textbf{Trigger}: Si verifica un errore prima del completamento della richiesta di riscrittura.
\end{itemize}

% __________________________________________________%


% UC79 - Richiesta Distant Writing
\hypertarget{uc79}{}
\subsubsection{UC 79: Richiesta Distant Writing}

\begin{figure}[H]
    \centering
    \includegraphics[width=0.6\textwidth]{richiesta_distant_writing.png}
    \caption{UC 79 – Richiesta Distant Writing}
\end{figure}

\begin{itemize}
    \item \textbf{Attore Principale}: Utente
    \item \textbf{Attore Secondario}: LLM
    \item \textbf{Precondizioni}: Una nota è aperta nell'editor, la finestra dedicata al Distant Writing è visualizzata e il servizio AI è disponibile.
    \item \textbf{Scenario Principale}:
    \begin{enumerate}
        \item L'utente inserisce un prompt nella finestra dedicata e richiede la generazione di un testo mediante Distant Writing.
        \item Il sistema acquisisce la richiesta dell'utente.
    \end{enumerate}
    \item \textbf{Postcondizioni}: La richiesta di Distant Writing è stata acquisita dal sistema.
    \item \textbf{Trigger}: L'utente desidera ottenere nuovo contenuto mediante Distant Writing.
\end{itemize}


% UC80 - Generazione Distant Writing
\hypertarget{uc80}{}
\subsubsection{UC 80: Generazione Distant Writing}

\begin{figure}[H]
    \centering
    \includegraphics[width=0.6\textwidth]{generazione_distant_writing.png}
    \caption{UC 80 – Generazione Distant Writing}
\end{figure}

\begin{itemize}
    \item \textbf{Attore Principale}: Utente
    \item \textbf{Attore Secondario}: LLM
    \item \textbf{Precondizioni}: La richiesta di Distant Writing è stata acquisita dal sistema.
    \item \textbf{Scenario Principale}:
    \begin{enumerate}
        \item Il sistema inoltra al servizio LLM la richiesta di Distant Writing.
        \item Il servizio LLM prende in carico la richiesta e avvia l'elaborazione.
    \end{enumerate}
    \item \textbf{Postcondizioni}: La richiesta di Distant Writing è in elaborazione presso il servizio LLM.
    \item \textbf{Trigger}: Il sistema ha acquisito la richiesta di Distant Writing.
\end{itemize}


% UC81 - Visualizzazione elaborazione Distant Writing
\hypertarget{uc81}{}
\subsubsection{UC 81: Visualizzazione elaborazione Distant Writing}

\begin{figure}[H]
    \centering
    \includegraphics[width=0.6\textwidth]{elaborazione_distant_writing.png}
    \caption{UC 81 – Visualizzazione elaborazione Distant Writing}
\end{figure}

\begin{itemize}
    \item \textbf{Attore Principale}: Utente
    \item \textbf{Precondizioni}: La richiesta di Distant Writing è in elaborazione presso il servizio LLM.
    \item \textbf{Scenario Principale}:
    \begin{enumerate}
        \item Il sistema visualizza una finestra modale contenente un indicatore di caricamento.
        \item Il sistema rende disponibile il comando ``Interrompi''.
    \end{enumerate}
    \item \textbf{Postcondizioni}: La finestra modale di elaborazione è visualizzata e la richiesta è ancora in corso.
    \item \textbf{Trigger}: Il servizio LLM ha avviato l'elaborazione della richiesta di Distant Writing.
\end{itemize}


% UC82 - Interruzione elaborazione Distant Writing
\hypertarget{uc82}{}
\subsubsection{UC 82: Interruzione elaborazione Distant Writing}

\begin{figure}[H]
    \centering
    \includegraphics[width=0.6\textwidth]{interruzione_distant_writing.png}
    \caption{UC 82 – Interruzione elaborazione Distant Writing}
\end{figure}

\begin{itemize}
    \item \textbf{Attore Principale}: Utente
    \item \textbf{Attore Secondario}: LLM
    \item \textbf{Precondizioni}: La finestra modale di elaborazione è visualizzata e la richiesta è ancora in corso.
    \item \textbf{Scenario Principale}:
    \begin{enumerate}
        \item L'utente attiva il comando ``Interrompi''.
        \item Il sistema richiede al servizio LLM l'interruzione dell'elaborazione.
        \item Il sistema chiude la finestra modale di elaborazione.
    \end{enumerate}
    \item \textbf{Postcondizioni}: L'elaborazione è interrotta e il contenuto della nota rimane invariato.
    \item \textbf{Trigger}: L'utente desidera interrompere la richiesta ancora in elaborazione.
\end{itemize}


% UC83 - Visualizzazione proposta Distant Writing
\hypertarget{uc83}{}
\subsubsection{UC 83: Visualizzazione proposta Distant Writing}

\begin{figure}[H]
    \centering
    \includegraphics[width=0.6\textwidth]{proposta_distant_writing.png}
    \caption{UC 83 – Visualizzazione proposta Distant Writing}
\end{figure}

\begin{itemize}
    \item \textbf{Attore Principale}: Utente
    \item \textbf{Precondizioni}: Il servizio LLM ha completato la generazione della proposta di Distant Writing.
    \item \textbf{Scenario Principale}:
    \begin{enumerate}
        \item Il sistema chiude la finestra modale di elaborazione.
        \item Il sistema visualizza il testo generato mediante Distant Writing in una finestra modale dedicata.
        \item Il sistema rende disponibili i comandi ``Accetta'', ``Rifiuta'' e ``Rigenera''.
    \end{enumerate}
    \item \textbf{Postcondizioni}: La proposta è visualizzata e il contenuto della nota rimane invariato.
    \item \textbf{Trigger}: Il servizio LLM ha restituito la proposta di Distant Writing.
\end{itemize}


% UC84 - Accettazione proposta Distant Writing
\hypertarget{uc84}{}
\subsubsection{UC 84: Accettazione proposta Distant Writing}

\begin{figure}[H]
    \centering
    \includegraphics[width=0.6\textwidth]{accettazione_distant_writing.png}
    \caption{UC 84 – Accettazione proposta Distant Writing}
\end{figure}

\begin{itemize}
    \item \textbf{Attore Principale}: Utente
    \item \textbf{Precondizioni}: La finestra modale contenente la proposta di Distant Writing è visualizzata.
    \item \textbf{Scenario Principale}:
    \begin{enumerate}
        \item L'utente attiva il comando ``Accetta''.
        \item Il sistema inserisce il testo generato mediante Distant Writing nella posizione corrente del cursore.
        \item Il sistema chiude la finestra modale della proposta.
    \end{enumerate}
    \item \textbf{Postcondizioni}: Il testo generato mediante Distant Writing è inserito nella nota e la finestra modale è chiusa.
    \item \textbf{Trigger}: L'utente desidera applicare la proposta visualizzata.
\end{itemize}


% UC85 - Rifiuto proposta Distant Writing
\hypertarget{uc85}{}
\subsubsection{UC 85: Rifiuto proposta Distant Writing}

\begin{figure}[H]
    \centering
    \includegraphics[width=0.6\textwidth]{rifiuto_distant_writing.png}
    \caption{UC 85 – Rifiuto proposta Distant Writing}
\end{figure}

\begin{itemize}
    \item \textbf{Attore Principale}: Utente
    \item \textbf{Precondizioni}: La finestra modale contenente la proposta di Distant Writing è visualizzata.
    \item \textbf{Scenario Principale}:
    \begin{enumerate}
        \item L'utente attiva il comando ``Rifiuta''.
        \item Il sistema scarta la proposta visualizzata.
        \item Il sistema chiude la finestra modale della proposta.
    \end{enumerate}
    \item \textbf{Postcondizioni}: La proposta è scartata, la finestra modale è chiusa e il contenuto della nota rimane invariato.
    \item \textbf{Trigger}: L'utente non desidera applicare la proposta visualizzata.
\end{itemize}


% UC86 - Rigenerazione proposta Distant Writing
\hypertarget{uc86}{}
\subsubsection{UC 86: Rigenerazione proposta Distant Writing}

\begin{figure}[H]
    \centering
    \includegraphics[width=0.6\textwidth]{rigenerazione_distant_writing.png}
    \caption{UC 86 – Rigenerazione proposta Distant Writing}
\end{figure}

\begin{itemize}
    \item \textbf{Attore Principale}: Utente
    \item \textbf{Attore Secondario}: LLM
    \item \textbf{Precondizioni}: La finestra modale contenente la proposta di Distant Writing è visualizzata.
    \item \textbf{Scenario Principale}:
    \begin{enumerate}
        \item L'utente attiva il comando ``Rigenera''.
        \item Il sistema scarta la proposta visualizzata.
        \item Il sistema inoltra al servizio LLM una nuova richiesta di Distant Writing.
        \item Il sistema visualizza nuovamente la finestra modale di elaborazione.
    \end{enumerate}
    \item \textbf{Postcondizioni}: Una nuova richiesta è in elaborazione e il contenuto della nota rimane invariato.
    \item \textbf{Trigger}: L'utente desidera ottenere una proposta differente.
\end{itemize}


% UC87 - Visualizzazione errore generazione Distant Writing
\hypertarget{uc87}{}
\subsubsection{UC 87: Visualizzazione errore generazione Distant Writing}

\begin{figure}[H]
    \centering
    \includegraphics[width=0.6\textwidth]{errore_distant_writing.png}
    \caption{UC 87 – Visualizzazione errore generazione Distant Writing}
\end{figure}

\begin{itemize}
    \item \textbf{Attore Principale}: Utente
    \item \textbf{Precondizioni}: È in corso una richiesta di Distant Writing.
    \item \textbf{Scenario Principale}:
    \begin{enumerate}
        \item Si verifica un errore durante la comunicazione con il servizio LLM o durante l'elaborazione della richiesta.
        \item Il sistema chiude la finestra modale di elaborazione.
        \item Il sistema visualizza un messaggio che segnala il mancato completamento della richiesta.
    \end{enumerate}
    \item \textbf{Postcondizioni}: L'utente è informato dell'errore e il contenuto della nota rimane invariato.
    \item \textbf{Trigger}: Si verifica un errore prima del completamento della richiesta di Distant Writing.
\end{itemize}

% __________________________________________________%


% UC88 - Richiesta analisi secondo il cappello bianco
\hypertarget{uc88}{}
\subsubsection{UC 88: Richiesta analisi secondo il cappello bianco}

\begin{figure}[H]
    \centering
    \includegraphics[width=0.6\textwidth]{richiesta_cappello_bianco.png}
    \caption{UC 88 – Richiesta analisi secondo il cappello bianco}
\end{figure}

\begin{itemize}
    \item \textbf{Attore Principale}: Utente
    \item \textbf{Attore Secondario}: LLM
    \item \textbf{Precondizioni}: Una porzione di testo è selezionata nell'editor e il servizio AI è disponibile.
    \item \textbf{Scenario Principale}:
    \begin{enumerate}
        \item L'utente richiede un'analisi del testo selezionato secondo il cappello bianco.
        \item Il sistema acquisisce la richiesta dell'utente.
    \end{enumerate}
    \item \textbf{Postcondizioni}: La richiesta di analisi secondo il cappello bianco è stata acquisita dal sistema.
    \item \textbf{Trigger}: L'utente desidera verificare dati, fatti e oggettività del testo selezionato.
\end{itemize}


% UC89 - Generazione analisi secondo il cappello bianco
\hypertarget{uc89}{}
\subsubsection{UC 89: Generazione analisi secondo il cappello bianco}

\begin{figure}[H]
    \centering
    \includegraphics[width=0.6\textwidth]{generazione_cappello_bianco.png}
    \caption{UC 89 – Generazione analisi secondo il cappello bianco}
\end{figure}

\begin{itemize}
    \item \textbf{Attore Principale}: Utente
    \item \textbf{Attore Secondario}: LLM
    \item \textbf{Precondizioni}: La richiesta di analisi secondo il cappello bianco è stata acquisita dal sistema.
    \item \textbf{Scenario Principale}:
    \begin{enumerate}
        \item Il sistema inoltra al servizio LLM la richiesta di analisi secondo il cappello bianco.
        \item Il servizio LLM prende in carico la richiesta e avvia l'elaborazione.
    \end{enumerate}
    \item \textbf{Postcondizioni}: La richiesta di analisi secondo il cappello bianco è in elaborazione presso il servizio LLM.
    \item \textbf{Trigger}: Il sistema ha acquisito la richiesta di analisi secondo il cappello bianco.
\end{itemize}


% UC90 - Visualizzazione elaborazione analisi secondo il cappello bianco
\hypertarget{uc90}{}
\subsubsection{UC 90: Visualizzazione elaborazione analisi secondo il cappello bianco}

\begin{figure}[H]
    \centering
    \includegraphics[width=0.6\textwidth]{elaborazione_cappello_bianco.png}
    \caption{UC 90 – Visualizzazione elaborazione analisi secondo il cappello bianco}
\end{figure}

\begin{itemize}
    \item \textbf{Attore Principale}: Utente
    \item \textbf{Precondizioni}: La richiesta di analisi secondo il cappello bianco è in elaborazione presso il servizio LLM.
    \item \textbf{Scenario Principale}:
    \begin{enumerate}
        \item Il sistema visualizza una finestra modale contenente un indicatore di caricamento.
        \item Il sistema rende disponibile il comando ``Interrompi''.
    \end{enumerate}
    \item \textbf{Postcondizioni}: La finestra modale di elaborazione è visualizzata e la richiesta è ancora in corso.
    \item \textbf{Trigger}: Il servizio LLM ha avviato l'elaborazione della richiesta di analisi secondo il cappello bianco.
\end{itemize}


% UC91 - Interruzione elaborazione analisi secondo il cappello bianco
\hypertarget{uc91}{}
\subsubsection{UC 91: Interruzione elaborazione analisi secondo il cappello bianco}

\begin{figure}[H]
    \centering
    \includegraphics[width=0.6\textwidth]{interruzione_cappello_bianco.png}
    \caption{UC 91 – Interruzione elaborazione analisi secondo il cappello bianco}
\end{figure}

\begin{itemize}
    \item \textbf{Attore Principale}: Utente
    \item \textbf{Attore Secondario}: LLM
    \item \textbf{Precondizioni}: La finestra modale di elaborazione è visualizzata e la richiesta è ancora in corso.
    \item \textbf{Scenario Principale}:
    \begin{enumerate}
        \item L'utente attiva il comando ``Interrompi''.
        \item Il sistema richiede al servizio LLM l'interruzione dell'elaborazione.
        \item Il sistema chiude la finestra modale di elaborazione.
    \end{enumerate}
    \item \textbf{Postcondizioni}: L'elaborazione è interrotta e il contenuto della nota rimane invariato.
    \item \textbf{Trigger}: L'utente desidera interrompere la richiesta ancora in elaborazione.
\end{itemize}


% UC92 - Visualizzazione proposta analisi secondo il cappello bianco
\hypertarget{uc92}{}
\subsubsection{UC 92: Visualizzazione proposta analisi secondo il cappello bianco}

\begin{figure}[H]
    \centering
    \includegraphics[width=0.6\textwidth]{proposta_cappello_bianco.png}
    \caption{UC 92 – Visualizzazione proposta analisi secondo il cappello bianco}
\end{figure}

\begin{itemize}
    \item \textbf{Attore Principale}: Utente
    \item \textbf{Precondizioni}: Il servizio LLM ha completato la generazione della proposta di analisi secondo il cappello bianco.
    \item \textbf{Scenario Principale}:
    \begin{enumerate}
        \item Il sistema chiude la finestra modale di elaborazione.
        \item Il sistema visualizza l'analisi secondo il cappello bianco in una finestra modale dedicata.
        \item Il sistema rende disponibili i comandi ``Accetta'', ``Rifiuta'' e ``Rigenera''.
    \end{enumerate}
    \item \textbf{Postcondizioni}: La proposta è visualizzata e il contenuto della nota rimane invariato.
    \item \textbf{Trigger}: Il servizio LLM ha restituito la proposta di analisi secondo il cappello bianco.
\end{itemize}


% UC93 - Accettazione proposta analisi secondo il cappello bianco
\hypertarget{uc93}{}
\subsubsection{UC 93: Accettazione proposta analisi secondo il cappello bianco}

\begin{figure}[H]
    \centering
    \includegraphics[width=0.6\textwidth]{accettazione_cappello_bianco.png}
    \caption{UC 93 – Accettazione proposta analisi secondo il cappello bianco}
\end{figure}

\begin{itemize}
    \item \textbf{Attore Principale}: Utente
    \item \textbf{Precondizioni}: La finestra modale contenente la proposta di analisi secondo il cappello bianco è visualizzata.
    \item \textbf{Scenario Principale}:
    \begin{enumerate}
        \item L'utente attiva il comando ``Accetta''.
        \item Il sistema sostituisce il testo inizialmente selezionato con l'analisi proposta secondo il cappello bianco.
        \item Il sistema chiude la finestra modale della proposta.
    \end{enumerate}
    \item \textbf{Postcondizioni}: Il testo inizialmente selezionato è sostituito con l'analisi secondo il cappello bianco accettata e la finestra modale è chiusa.
    \item \textbf{Trigger}: L'utente desidera applicare la proposta visualizzata.
\end{itemize}


% UC94 - Rifiuto proposta analisi secondo il cappello bianco
\hypertarget{uc94}{}
\subsubsection{UC 94: Rifiuto proposta analisi secondo il cappello bianco}

\begin{figure}[H]
    \centering
    \includegraphics[width=0.6\textwidth]{rifiuto_cappello_bianco.png}
    \caption{UC 94 – Rifiuto proposta analisi secondo il cappello bianco}
\end{figure}

\begin{itemize}
    \item \textbf{Attore Principale}: Utente
    \item \textbf{Precondizioni}: La finestra modale contenente la proposta di analisi secondo il cappello bianco è visualizzata.
    \item \textbf{Scenario Principale}:
    \begin{enumerate}
        \item L'utente attiva il comando ``Rifiuta''.
        \item Il sistema scarta la proposta visualizzata.
        \item Il sistema chiude la finestra modale della proposta.
    \end{enumerate}
    \item \textbf{Postcondizioni}: La proposta è scartata, la finestra modale è chiusa e il contenuto della nota rimane invariato.
    \item \textbf{Trigger}: L'utente non desidera applicare la proposta visualizzata.
\end{itemize}


% UC95 - Rigenerazione proposta analisi secondo il cappello bianco
\hypertarget{uc95}{}
\subsubsection{UC 95: Rigenerazione proposta analisi secondo il cappello bianco}

\begin{figure}[H]
    \centering
    \includegraphics[width=0.6\textwidth]{rigenerazione_cappello_bianco.png}
    \caption{UC 95 – Rigenerazione proposta analisi secondo il cappello bianco}
\end{figure}

\begin{itemize}
    \item \textbf{Attore Principale}: Utente
    \item \textbf{Attore Secondario}: LLM
    \item \textbf{Precondizioni}: La finestra modale contenente la proposta di analisi secondo il cappello bianco è visualizzata.
    \item \textbf{Scenario Principale}:
    \begin{enumerate}
        \item L'utente attiva il comando ``Rigenera''.
        \item Il sistema scarta la proposta visualizzata.
        \item Il sistema inoltra al servizio LLM una nuova richiesta di analisi secondo il cappello bianco.
        \item Il sistema visualizza nuovamente la finestra modale di elaborazione.
    \end{enumerate}
    \item \textbf{Postcondizioni}: Una nuova richiesta è in elaborazione e il contenuto della nota rimane invariato.
    \item \textbf{Trigger}: L'utente desidera ottenere una proposta differente.
\end{itemize}


% UC96 - Visualizzazione errore generazione analisi secondo il cappello bianco
\hypertarget{uc96}{}
\subsubsection{UC 96: Visualizzazione errore generazione analisi secondo il cappello bianco}

\begin{figure}[H]
    \centering
    \includegraphics[width=0.6\textwidth]{errore_cappello_bianco.png}
    \caption{UC 96 – Visualizzazione errore generazione analisi secondo il cappello bianco}
\end{figure}

\begin{itemize}
    \item \textbf{Attore Principale}: Utente
    \item \textbf{Precondizioni}: È in corso una richiesta di analisi secondo il cappello bianco.
    \item \textbf{Scenario Principale}:
    \begin{enumerate}
        \item Si verifica un errore durante la comunicazione con il servizio LLM o durante l'elaborazione della richiesta.
        \item Il sistema chiude la finestra modale di elaborazione.
        \item Il sistema visualizza un messaggio che segnala il mancato completamento della richiesta.
    \end{enumerate}
    \item \textbf{Postcondizioni}: L'utente è informato dell'errore e il contenuto della nota rimane invariato.
    \item \textbf{Trigger}: Si verifica un errore prima del completamento della richiesta di analisi secondo il cappello bianco.
\end{itemize}

% __________________________________________________%


% UC97 - Richiesta analisi secondo il cappello rosso
\hypertarget{uc97}{}
\subsubsection{UC 97: Richiesta analisi secondo il cappello rosso}

\begin{figure}[H]
    \centering
    \includegraphics[width=0.6\textwidth]{richiesta_cappello_rosso.png}
    \caption{UC 97 – Richiesta analisi secondo il cappello rosso}
\end{figure}

\begin{itemize}
    \item \textbf{Attore Principale}: Utente
    \item \textbf{Attore Secondario}: LLM
    \item \textbf{Precondizioni}: Una porzione di testo è selezionata nell'editor e il servizio AI è disponibile.
    \item \textbf{Scenario Principale}:
    \begin{enumerate}
        \item L'utente richiede un'analisi del testo selezionato secondo il cappello rosso.
        \item Il sistema acquisisce la richiesta dell'utente.
    \end{enumerate}
    \item \textbf{Postcondizioni}: La richiesta di analisi secondo il cappello rosso è stata acquisita dal sistema.
    \item \textbf{Trigger}: L'utente desidera valutare l'impatto emotivo del testo selezionato.
\end{itemize}


% UC98 - Generazione analisi secondo il cappello rosso
\hypertarget{uc98}{}
\subsubsection{UC 98: Generazione analisi secondo il cappello rosso}

\begin{figure}[H]
    \centering
    \includegraphics[width=0.6\textwidth]{generazione_cappello_rosso.png}
    \caption{UC 98 – Generazione analisi secondo il cappello rosso}
\end{figure}

\begin{itemize}
    \item \textbf{Attore Principale}: Utente
    \item \textbf{Attore Secondario}: LLM
    \item \textbf{Precondizioni}: La richiesta di analisi secondo il cappello rosso è stata acquisita dal sistema.
    \item \textbf{Scenario Principale}:
    \begin{enumerate}
        \item Il sistema inoltra al servizio LLM la richiesta di analisi secondo il cappello rosso.
        \item Il servizio LLM prende in carico la richiesta e avvia l'elaborazione.
    \end{enumerate}
    \item \textbf{Postcondizioni}: La richiesta di analisi secondo il cappello rosso è in elaborazione presso il servizio LLM.
    \item \textbf{Trigger}: Il sistema ha acquisito la richiesta di analisi secondo il cappello rosso.
\end{itemize}


% UC99 - Visualizzazione elaborazione analisi secondo il cappello rosso
\hypertarget{uc99}{}
\subsubsection{UC 99: Visualizzazione elaborazione analisi secondo il cappello rosso}

\begin{figure}[H]
    \centering
    \includegraphics[width=0.6\textwidth]{elaborazione_cappello_rosso.png}
    \caption{UC 99 – Visualizzazione elaborazione analisi secondo il cappello rosso}
\end{figure}

\begin{itemize}
    \item \textbf{Attore Principale}: Utente
    \item \textbf{Precondizioni}: La richiesta di analisi secondo il cappello rosso è in elaborazione presso il servizio LLM.
    \item \textbf{Scenario Principale}:
    \begin{enumerate}
        \item Il sistema visualizza una finestra modale contenente un indicatore di caricamento.
        \item Il sistema rende disponibile il comando ``Interrompi''.
    \end{enumerate}
    \item \textbf{Postcondizioni}: La finestra modale di elaborazione è visualizzata e la richiesta è ancora in corso.
    \item \textbf{Trigger}: Il servizio LLM ha avviato l'elaborazione della richiesta di analisi secondo il cappello rosso.
\end{itemize}


% UC100 - Interruzione elaborazione analisi secondo il cappello rosso
\hypertarget{uc100}{}
\subsubsection{UC 100: Interruzione elaborazione analisi secondo il cappello rosso}

\begin{figure}[H]
    \centering
    \includegraphics[width=0.6\textwidth]{interruzione_cappello_rosso.png}
    \caption{UC 100 – Interruzione elaborazione analisi secondo il cappello rosso}
\end{figure}

\begin{itemize}
    \item \textbf{Attore Principale}: Utente
    \item \textbf{Attore Secondario}: LLM
    \item \textbf{Precondizioni}: La finestra modale di elaborazione è visualizzata e la richiesta è ancora in corso.
    \item \textbf{Scenario Principale}:
    \begin{enumerate}
        \item L'utente attiva il comando ``Interrompi''.
        \item Il sistema richiede al servizio LLM l'interruzione dell'elaborazione.
        \item Il sistema chiude la finestra modale di elaborazione.
    \end{enumerate}
    \item \textbf{Postcondizioni}: L'elaborazione è interrotta e il contenuto della nota rimane invariato.
    \item \textbf{Trigger}: L'utente desidera interrompere la richiesta ancora in elaborazione.
\end{itemize}


% UC101 - Visualizzazione proposta analisi secondo il cappello rosso
\hypertarget{uc101}{}
\subsubsection{UC 101: Visualizzazione proposta analisi secondo il cappello rosso}

\begin{figure}[H]
    \centering
    \includegraphics[width=0.6\textwidth]{proposta_cappello_rosso.png}
    \caption{UC 101 – Visualizzazione proposta analisi secondo il cappello rosso}
\end{figure}

\begin{itemize}
    \item \textbf{Attore Principale}: Utente
    \item \textbf{Precondizioni}: Il servizio LLM ha completato la generazione della proposta di analisi secondo il cappello rosso.
    \item \textbf{Scenario Principale}:
    \begin{enumerate}
        \item Il sistema chiude la finestra modale di elaborazione.
        \item Il sistema visualizza l'analisi secondo il cappello rosso in una finestra modale dedicata.
        \item Il sistema rende disponibili i comandi ``Accetta'', ``Rifiuta'' e ``Rigenera''.
    \end{enumerate}
    \item \textbf{Postcondizioni}: La proposta è visualizzata e il contenuto della nota rimane invariato.
    \item \textbf{Trigger}: Il servizio LLM ha restituito la proposta di analisi secondo il cappello rosso.
\end{itemize}


% UC102 - Accettazione proposta analisi secondo il cappello rosso
\hypertarget{uc102}{}
\subsubsection{UC 102: Accettazione proposta analisi secondo il cappello rosso}

\begin{figure}[H]
    \centering
    \includegraphics[width=0.6\textwidth]{accettazione_cappello_rosso.png}
    \caption{UC 102 – Accettazione proposta analisi secondo il cappello rosso}
\end{figure}

\begin{itemize}
    \item \textbf{Attore Principale}: Utente
    \item \textbf{Precondizioni}: La finestra modale contenente la proposta di analisi secondo il cappello rosso è visualizzata.
    \item \textbf{Scenario Principale}:
    \begin{enumerate}
        \item L'utente attiva il comando ``Accetta''.
        \item Il sistema sostituisce il testo inizialmente selezionato con l'analisi proposta secondo il cappello rosso.
        \item Il sistema chiude la finestra modale della proposta.
    \end{enumerate}
    \item \textbf{Postcondizioni}: Il testo inizialmente selezionato è sostituito con l'analisi secondo il cappello rosso accettata e la finestra modale è chiusa.
    \item \textbf{Trigger}: L'utente desidera applicare la proposta visualizzata.
\end{itemize}


% UC103 - Rifiuto proposta analisi secondo il cappello rosso
\hypertarget{uc103}{}
\subsubsection{UC 103: Rifiuto proposta analisi secondo il cappello rosso}

\begin{figure}[H]
    \centering
    \includegraphics[width=0.6\textwidth]{rifiuto_cappello_rosso.png}
    \caption{UC 103 – Rifiuto proposta analisi secondo il cappello rosso}
\end{figure}

\begin{itemize}
    \item \textbf{Attore Principale}: Utente
    \item \textbf{Precondizioni}: La finestra modale contenente la proposta di analisi secondo il cappello rosso è visualizzata.
    \item \textbf{Scenario Principale}:
    \begin{enumerate}
        \item L'utente attiva il comando ``Rifiuta''.
        \item Il sistema scarta la proposta visualizzata.
        \item Il sistema chiude la finestra modale della proposta.
    \end{enumerate}
    \item \textbf{Postcondizioni}: La proposta è scartata, la finestra modale è chiusa e il contenuto della nota rimane invariato.
    \item \textbf{Trigger}: L'utente non desidera applicare la proposta visualizzata.
\end{itemize}


% UC104 - Rigenerazione proposta analisi secondo il cappello rosso
\hypertarget{uc104}{}
\subsubsection{UC 104: Rigenerazione proposta analisi secondo il cappello rosso}

\begin{figure}[H]
    \centering
    \includegraphics[width=0.6\textwidth]{rigenerazione_cappello_rosso.png}
    \caption{UC 104 – Rigenerazione proposta analisi secondo il cappello rosso}
\end{figure}

\begin{itemize}
    \item \textbf{Attore Principale}: Utente
    \item \textbf{Attore Secondario}: LLM
    \item \textbf{Precondizioni}: La finestra modale contenente la proposta di analisi secondo il cappello rosso è visualizzata.
    \item \textbf{Scenario Principale}:
    \begin{enumerate}
        \item L'utente attiva il comando ``Rigenera''.
        \item Il sistema scarta la proposta visualizzata.
        \item Il sistema inoltra al servizio LLM una nuova richiesta di analisi secondo il cappello rosso.
        \item Il sistema visualizza nuovamente la finestra modale di elaborazione.
    \end{enumerate}
    \item \textbf{Postcondizioni}: Una nuova richiesta è in elaborazione e il contenuto della nota rimane invariato.
    \item \textbf{Trigger}: L'utente desidera ottenere una proposta differente.
\end{itemize}


% UC105 - Visualizzazione errore generazione analisi secondo il cappello rosso
\hypertarget{uc105}{}
\subsubsection{UC 105: Visualizzazione errore generazione analisi secondo il cappello rosso}

\begin{figure}[H]
    \centering
    \includegraphics[width=0.6\textwidth]{errore_cappello_rosso.png}
    \caption{UC 105 – Visualizzazione errore generazione analisi secondo il cappello rosso}
\end{figure}

\begin{itemize}
    \item \textbf{Attore Principale}: Utente
    \item \textbf{Precondizioni}: È in corso una richiesta di analisi secondo il cappello rosso.
    \item \textbf{Scenario Principale}:
    \begin{enumerate}
        \item Si verifica un errore durante la comunicazione con il servizio LLM o durante l'elaborazione della richiesta.
        \item Il sistema chiude la finestra modale di elaborazione.
        \item Il sistema visualizza un messaggio che segnala il mancato completamento della richiesta.
    \end{enumerate}
    \item \textbf{Postcondizioni}: L'utente è informato dell'errore e il contenuto della nota rimane invariato.
    \item \textbf{Trigger}: Si verifica un errore prima del completamento della richiesta di analisi secondo il cappello rosso.
\end{itemize}

% __________________________________________________%


% UC106 - Richiesta analisi secondo il cappello nero
\hypertarget{uc106}{}
\subsubsection{UC 106: Richiesta analisi secondo il cappello nero}

\begin{figure}[H]
    \centering
    \includegraphics[width=0.6\textwidth]{richiesta_cappello_nero.png}
    \caption{UC 106 – Richiesta analisi secondo il cappello nero}
\end{figure}

\begin{itemize}
    \item \textbf{Attore Principale}: Utente
    \item \textbf{Attore Secondario}: LLM
    \item \textbf{Precondizioni}: Una porzione di testo è selezionata nell'editor e il servizio AI è disponibile.
    \item \textbf{Scenario Principale}:
    \begin{enumerate}
        \item L'utente richiede un'analisi del testo selezionato secondo il cappello nero.
        \item Il sistema acquisisce la richiesta dell'utente.
    \end{enumerate}
    \item \textbf{Postcondizioni}: La richiesta di analisi secondo il cappello nero è stata acquisita dal sistema.
    \item \textbf{Trigger}: L'utente desidera individuare rischi, criticità e punti deboli del testo selezionato.
\end{itemize}


% UC107 - Generazione analisi secondo il cappello nero
\hypertarget{uc107}{}
\subsubsection{UC 107: Generazione analisi secondo il cappello nero}

\begin{figure}[H]
    \centering
    \includegraphics[width=0.6\textwidth]{generazione_cappello_nero.png}
    \caption{UC 107 – Generazione analisi secondo il cappello nero}
\end{figure}

\begin{itemize}
    \item \textbf{Attore Principale}: Utente
    \item \textbf{Attore Secondario}: LLM
    \item \textbf{Precondizioni}: La richiesta di analisi secondo il cappello nero è stata acquisita dal sistema.
    \item \textbf{Scenario Principale}:
    \begin{enumerate}
        \item Il sistema inoltra al servizio LLM la richiesta di analisi secondo il cappello nero.
        \item Il servizio LLM prende in carico la richiesta e avvia l'elaborazione.
    \end{enumerate}
    \item \textbf{Postcondizioni}: La richiesta di analisi secondo il cappello nero è in elaborazione presso il servizio LLM.
    \item \textbf{Trigger}: Il sistema ha acquisito la richiesta di analisi secondo il cappello nero.
\end{itemize}


% UC108 - Visualizzazione elaborazione analisi secondo il cappello nero
\hypertarget{uc108}{}
\subsubsection{UC 108: Visualizzazione elaborazione analisi secondo il cappello nero}

\begin{figure}[H]
    \centering
    \includegraphics[width=0.6\textwidth]{elaborazione_cappello_nero.png}
    \caption{UC 108 – Visualizzazione elaborazione analisi secondo il cappello nero}
\end{figure}

\begin{itemize}
    \item \textbf{Attore Principale}: Utente
    \item \textbf{Precondizioni}: La richiesta di analisi secondo il cappello nero è in elaborazione presso il servizio LLM.
    \item \textbf{Scenario Principale}:
    \begin{enumerate}
        \item Il sistema visualizza una finestra modale contenente un indicatore di caricamento.
        \item Il sistema rende disponibile il comando ``Interrompi''.
    \end{enumerate}
    \item \textbf{Postcondizioni}: La finestra modale di elaborazione è visualizzata e la richiesta è ancora in corso.
    \item \textbf{Trigger}: Il servizio LLM ha avviato l'elaborazione della richiesta di analisi secondo il cappello nero.
\end{itemize}


% UC109 - Interruzione elaborazione analisi secondo il cappello nero
\hypertarget{uc109}{}
\subsubsection{UC 109: Interruzione elaborazione analisi secondo il cappello nero}

\begin{figure}[H]
    \centering
    \includegraphics[width=0.6\textwidth]{interruzione_cappello_nero.png}
    \caption{UC 109 – Interruzione elaborazione analisi secondo il cappello nero}
\end{figure}

\begin{itemize}
    \item \textbf{Attore Principale}: Utente
    \item \textbf{Attore Secondario}: LLM
    \item \textbf{Precondizioni}: La finestra modale di elaborazione è visualizzata e la richiesta è ancora in corso.
    \item \textbf{Scenario Principale}:
    \begin{enumerate}
        \item L'utente attiva il comando ``Interrompi''.
        \item Il sistema richiede al servizio LLM l'interruzione dell'elaborazione.
        \item Il sistema chiude la finestra modale di elaborazione.
    \end{enumerate}
    \item \textbf{Postcondizioni}: L'elaborazione è interrotta e il contenuto della nota rimane invariato.
    \item \textbf{Trigger}: L'utente desidera interrompere la richiesta ancora in elaborazione.
\end{itemize}


% UC110 - Visualizzazione proposta analisi secondo il cappello nero
\hypertarget{uc110}{}
\subsubsection{UC 110: Visualizzazione proposta analisi secondo il cappello nero}

\begin{figure}[H]
    \centering
    \includegraphics[width=0.6\textwidth]{proposta_cappello_nero.png}
    \caption{UC 110 – Visualizzazione proposta analisi secondo il cappello nero}
\end{figure}

\begin{itemize}
    \item \textbf{Attore Principale}: Utente
    \item \textbf{Precondizioni}: Il servizio LLM ha completato la generazione della proposta di analisi secondo il cappello nero.
    \item \textbf{Scenario Principale}:
    \begin{enumerate}
        \item Il sistema chiude la finestra modale di elaborazione.
        \item Il sistema visualizza l'analisi secondo il cappello nero in una finestra modale dedicata.
        \item Il sistema rende disponibili i comandi ``Accetta'', ``Rifiuta'' e ``Rigenera''.
    \end{enumerate}
    \item \textbf{Postcondizioni}: La proposta è visualizzata e il contenuto della nota rimane invariato.
    \item \textbf{Trigger}: Il servizio LLM ha restituito la proposta di analisi secondo il cappello nero.
\end{itemize}


% UC111 - Accettazione proposta analisi secondo il cappello nero
\hypertarget{uc111}{}
\subsubsection{UC 111: Accettazione proposta analisi secondo il cappello nero}

\begin{figure}[H]
    \centering
    \includegraphics[width=0.6\textwidth]{accettazione_cappello_nero.png}
    \caption{UC 111 – Accettazione proposta analisi secondo il cappello nero}
\end{figure}

\begin{itemize}
    \item \textbf{Attore Principale}: Utente
    \item \textbf{Precondizioni}: La finestra modale contenente la proposta di analisi secondo il cappello nero è visualizzata.
    \item \textbf{Scenario Principale}:
    \begin{enumerate}
        \item L'utente attiva il comando ``Accetta''.
        \item Il sistema sostituisce il testo inizialmente selezionato con l'analisi proposta secondo il cappello nero.
        \item Il sistema chiude la finestra modale della proposta.
    \end{enumerate}
    \item \textbf{Postcondizioni}: Il testo inizialmente selezionato è sostituito con l'analisi secondo il cappello nero accettata e la finestra modale è chiusa.
    \item \textbf{Trigger}: L'utente desidera applicare la proposta visualizzata.
\end{itemize}


% UC112 - Rifiuto proposta analisi secondo il cappello nero
\hypertarget{uc112}{}
\subsubsection{UC 112: Rifiuto proposta analisi secondo il cappello nero}

\begin{figure}[H]
    \centering
    \includegraphics[width=0.6\textwidth]{rifiuto_cappello_nero.png}
    \caption{UC 112 – Rifiuto proposta analisi secondo il cappello nero}
\end{figure}

\begin{itemize}
    \item \textbf{Attore Principale}: Utente
    \item \textbf{Precondizioni}: La finestra modale contenente la proposta di analisi secondo il cappello nero è visualizzata.
    \item \textbf{Scenario Principale}:
    \begin{enumerate}
        \item L'utente attiva il comando ``Rifiuta''.
        \item Il sistema scarta la proposta visualizzata.
        \item Il sistema chiude la finestra modale della proposta.
    \end{enumerate}
    \item \textbf{Postcondizioni}: La proposta è scartata, la finestra modale è chiusa e il contenuto della nota rimane invariato.
    \item \textbf{Trigger}: L'utente non desidera applicare la proposta visualizzata.
\end{itemize}


% UC113 - Rigenerazione proposta analisi secondo il cappello nero
\hypertarget{uc113}{}
\subsubsection{UC 113: Rigenerazione proposta analisi secondo il cappello nero}

\begin{figure}[H]
    \centering
    \includegraphics[width=0.6\textwidth]{rigenerazione_cappello_nero.png}
    \caption{UC 113 – Rigenerazione proposta analisi secondo il cappello nero}
\end{figure}

\begin{itemize}
    \item \textbf{Attore Principale}: Utente
    \item \textbf{Attore Secondario}: LLM
    \item \textbf{Precondizioni}: La finestra modale contenente la proposta di analisi secondo il cappello nero è visualizzata.
    \item \textbf{Scenario Principale}:
    \begin{enumerate}
        \item L'utente attiva il comando ``Rigenera''.
        \item Il sistema scarta la proposta visualizzata.
        \item Il sistema inoltra al servizio LLM una nuova richiesta di analisi secondo il cappello nero.
        \item Il sistema visualizza nuovamente la finestra modale di elaborazione.
    \end{enumerate}
    \item \textbf{Postcondizioni}: Una nuova richiesta è in elaborazione e il contenuto della nota rimane invariato.
    \item \textbf{Trigger}: L'utente desidera ottenere una proposta differente.
\end{itemize}


% UC114 - Visualizzazione errore generazione analisi secondo il cappello nero
\hypertarget{uc114}{}
\subsubsection{UC 114: Visualizzazione errore generazione analisi secondo il cappello nero}

\begin{figure}[H]
    \centering
    \includegraphics[width=0.6\textwidth]{errore_cappello_nero.png}
    \caption{UC 114 – Visualizzazione errore generazione analisi secondo il cappello nero}
\end{figure}

\begin{itemize}
    \item \textbf{Attore Principale}: Utente
    \item \textbf{Precondizioni}: È in corso una richiesta di analisi secondo il cappello nero.
    \item \textbf{Scenario Principale}:
    \begin{enumerate}
        \item Si verifica un errore durante la comunicazione con il servizio LLM o durante l'elaborazione della richiesta.
        \item Il sistema chiude la finestra modale di elaborazione.
        \item Il sistema visualizza un messaggio che segnala il mancato completamento della richiesta.
    \end{enumerate}
    \item \textbf{Postcondizioni}: L'utente è informato dell'errore e il contenuto della nota rimane invariato.
    \item \textbf{Trigger}: Si verifica un errore prima del completamento della richiesta di analisi secondo il cappello nero.
\end{itemize}

% __________________________________________________%


% UC115 - Richiesta analisi secondo il cappello giallo
\hypertarget{uc115}{}
\subsubsection{UC 115: Richiesta analisi secondo il cappello giallo}

\begin{figure}[H]
    \centering
    \includegraphics[width=0.6\textwidth]{richiesta_cappello_giallo.png}
    \caption{UC 115 – Richiesta analisi secondo il cappello giallo}
\end{figure}

\begin{itemize}
    \item \textbf{Attore Principale}: Utente
    \item \textbf{Attore Secondario}: LLM
    \item \textbf{Precondizioni}: Una porzione di testo è selezionata nell'editor e il servizio AI è disponibile.
    \item \textbf{Scenario Principale}:
    \begin{enumerate}
        \item L'utente richiede un'analisi del testo selezionato secondo il cappello giallo.
        \item Il sistema acquisisce la richiesta dell'utente.
    \end{enumerate}
    \item \textbf{Postcondizioni}: La richiesta di analisi secondo il cappello giallo è stata acquisita dal sistema.
    \item \textbf{Trigger}: L'utente desidera individuare vantaggi e opportunità del testo selezionato.
\end{itemize}


% UC116 - Generazione analisi secondo il cappello giallo
\hypertarget{uc116}{}
\subsubsection{UC 116: Generazione analisi secondo il cappello giallo}

\begin{figure}[H]
    \centering
    \includegraphics[width=0.6\textwidth]{generazione_cappello_giallo.png}
    \caption{UC 116 – Generazione analisi secondo il cappello giallo}
\end{figure}

\begin{itemize}
    \item \textbf{Attore Principale}: Utente
    \item \textbf{Attore Secondario}: LLM
    \item \textbf{Precondizioni}: La richiesta di analisi secondo il cappello giallo è stata acquisita dal sistema.
    \item \textbf{Scenario Principale}:
    \begin{enumerate}
        \item Il sistema inoltra al servizio LLM la richiesta di analisi secondo il cappello giallo.
        \item Il servizio LLM prende in carico la richiesta e avvia l'elaborazione.
    \end{enumerate}
    \item \textbf{Postcondizioni}: La richiesta di analisi secondo il cappello giallo è in elaborazione presso il servizio LLM.
    \item \textbf{Trigger}: Il sistema ha acquisito la richiesta di analisi secondo il cappello giallo.
\end{itemize}


% UC117 - Visualizzazione elaborazione analisi secondo il cappello giallo
\hypertarget{uc117}{}
\subsubsection{UC 117: Visualizzazione elaborazione analisi secondo il cappello giallo}

\begin{figure}[H]
    \centering
    \includegraphics[width=0.6\textwidth]{elaborazione_cappello_giallo.png}
    \caption{UC 117 – Visualizzazione elaborazione analisi secondo il cappello giallo}
\end{figure}

\begin{itemize}
    \item \textbf{Attore Principale}: Utente
    \item \textbf{Precondizioni}: La richiesta di analisi secondo il cappello giallo è in elaborazione presso il servizio LLM.
    \item \textbf{Scenario Principale}:
    \begin{enumerate}
        \item Il sistema visualizza una finestra modale contenente un indicatore di caricamento.
        \item Il sistema rende disponibile il comando ``Interrompi''.
    \end{enumerate}
    \item \textbf{Postcondizioni}: La finestra modale di elaborazione è visualizzata e la richiesta è ancora in corso.
    \item \textbf{Trigger}: Il servizio LLM ha avviato l'elaborazione della richiesta di analisi secondo il cappello giallo.
\end{itemize}


% UC118 - Interruzione elaborazione analisi secondo il cappello giallo
\hypertarget{uc118}{}
\subsubsection{UC 118: Interruzione elaborazione analisi secondo il cappello giallo}

\begin{figure}[H]
    \centering
    \includegraphics[width=0.6\textwidth]{interruzione_cappello_giallo.png}
    \caption{UC 118 – Interruzione elaborazione analisi secondo il cappello giallo}
\end{figure}

\begin{itemize}
    \item \textbf{Attore Principale}: Utente
    \item \textbf{Attore Secondario}: LLM
    \item \textbf{Precondizioni}: La finestra modale di elaborazione è visualizzata e la richiesta è ancora in corso.
    \item \textbf{Scenario Principale}:
    \begin{enumerate}
        \item L'utente attiva il comando ``Interrompi''.
        \item Il sistema richiede al servizio LLM l'interruzione dell'elaborazione.
        \item Il sistema chiude la finestra modale di elaborazione.
    \end{enumerate}
    \item \textbf{Postcondizioni}: L'elaborazione è interrotta e il contenuto della nota rimane invariato.
    \item \textbf{Trigger}: L'utente desidera interrompere la richiesta ancora in elaborazione.
\end{itemize}


% UC119 - Visualizzazione proposta analisi secondo il cappello giallo
\hypertarget{uc119}{}
\subsubsection{UC 119: Visualizzazione proposta analisi secondo il cappello giallo}

\begin{figure}[H]
    \centering
    \includegraphics[width=0.6\textwidth]{proposta_cappello_giallo.png}
    \caption{UC 119 – Visualizzazione proposta analisi secondo il cappello giallo}
\end{figure}

\begin{itemize}
    \item \textbf{Attore Principale}: Utente
    \item \textbf{Precondizioni}: Il servizio LLM ha completato la generazione della proposta di analisi secondo il cappello giallo.
    \item \textbf{Scenario Principale}:
    \begin{enumerate}
        \item Il sistema chiude la finestra modale di elaborazione.
        \item Il sistema visualizza l'analisi secondo il cappello giallo in una finestra modale dedicata.
        \item Il sistema rende disponibili i comandi ``Accetta'', ``Rifiuta'' e ``Rigenera''.
    \end{enumerate}
    \item \textbf{Postcondizioni}: La proposta è visualizzata e il contenuto della nota rimane invariato.
    \item \textbf{Trigger}: Il servizio LLM ha restituito la proposta di analisi secondo il cappello giallo.
\end{itemize}


% UC120 - Accettazione proposta analisi secondo il cappello giallo
\hypertarget{uc120}{}
\subsubsection{UC 120: Accettazione proposta analisi secondo il cappello giallo}

\begin{figure}[H]
    \centering
    \includegraphics[width=0.6\textwidth]{accettazione_cappello_giallo.png}
    \caption{UC 120 – Accettazione proposta analisi secondo il cappello giallo}
\end{figure}

\begin{itemize}
    \item \textbf{Attore Principale}: Utente
    \item \textbf{Precondizioni}: La finestra modale contenente la proposta di analisi secondo il cappello giallo è visualizzata.
    \item \textbf{Scenario Principale}:
    \begin{enumerate}
        \item L'utente attiva il comando ``Accetta''.
        \item Il sistema sostituisce il testo inizialmente selezionato con l'analisi proposta secondo il cappello giallo.
        \item Il sistema chiude la finestra modale della proposta.
    \end{enumerate}
    \item \textbf{Postcondizioni}: Il testo inizialmente selezionato è sostituito con l'analisi secondo il cappello giallo accettata e la finestra modale è chiusa.
    \item \textbf{Trigger}: L'utente desidera applicare la proposta visualizzata.
\end{itemize}


% UC121 - Rifiuto proposta analisi secondo il cappello giallo
\hypertarget{uc121}{}
\subsubsection{UC 121: Rifiuto proposta analisi secondo il cappello giallo}

\begin{figure}[H]
    \centering
    \includegraphics[width=0.6\textwidth]{rifiuto_cappello_giallo.png}
    \caption{UC 121 – Rifiuto proposta analisi secondo il cappello giallo}
\end{figure}

\begin{itemize}
    \item \textbf{Attore Principale}: Utente
    \item \textbf{Precondizioni}: La finestra modale contenente la proposta di analisi secondo il cappello giallo è visualizzata.
    \item \textbf{Scenario Principale}:
    \begin{enumerate}
        \item L'utente attiva il comando ``Rifiuta''.
        \item Il sistema scarta la proposta visualizzata.
        \item Il sistema chiude la finestra modale della proposta.
    \end{enumerate}
    \item \textbf{Postcondizioni}: La proposta è scartata, la finestra modale è chiusa e il contenuto della nota rimane invariato.
    \item \textbf{Trigger}: L'utente non desidera applicare la proposta visualizzata.
\end{itemize}


% UC122 - Rigenerazione proposta analisi secondo il cappello giallo
\hypertarget{uc122}{}
\subsubsection{UC 122: Rigenerazione proposta analisi secondo il cappello giallo}

\begin{figure}[H]
    \centering
    \includegraphics[width=0.6\textwidth]{rigenerazione_cappello_giallo.png}
    \caption{UC 122 – Rigenerazione proposta analisi secondo il cappello giallo}
\end{figure}

\begin{itemize}
    \item \textbf{Attore Principale}: Utente
    \item \textbf{Attore Secondario}: LLM
    \item \textbf{Precondizioni}: La finestra modale contenente la proposta di analisi secondo il cappello giallo è visualizzata.
    \item \textbf{Scenario Principale}:
    \begin{enumerate}
        \item L'utente attiva il comando ``Rigenera''.
        \item Il sistema scarta la proposta visualizzata.
        \item Il sistema inoltra al servizio LLM una nuova richiesta di analisi secondo il cappello giallo.
        \item Il sistema visualizza nuovamente la finestra modale di elaborazione.
    \end{enumerate}
    \item \textbf{Postcondizioni}: Una nuova richiesta è in elaborazione e il contenuto della nota rimane invariato.
    \item \textbf{Trigger}: L'utente desidera ottenere una proposta differente.
\end{itemize}


% UC123 - Visualizzazione errore generazione analisi secondo il cappello giallo
\hypertarget{uc123}{}
\subsubsection{UC 123: Visualizzazione errore generazione analisi secondo il cappello giallo}

\begin{figure}[H]
    \centering
    \includegraphics[width=0.6\textwidth]{errore_cappello_giallo.png}
    \caption{UC 123 – Visualizzazione errore generazione analisi secondo il cappello giallo}
\end{figure}

\begin{itemize}
    \item \textbf{Attore Principale}: Utente
    \item \textbf{Precondizioni}: È in corso una richiesta di analisi secondo il cappello giallo.
    \item \textbf{Scenario Principale}:
    \begin{enumerate}
        \item Si verifica un errore durante la comunicazione con il servizio LLM o durante l'elaborazione della richiesta.
        \item Il sistema chiude la finestra modale di elaborazione.
        \item Il sistema visualizza un messaggio che segnala il mancato completamento della richiesta.
    \end{enumerate}
    \item \textbf{Postcondizioni}: L'utente è informato dell'errore e il contenuto della nota rimane invariato.
    \item \textbf{Trigger}: Si verifica un errore prima del completamento della richiesta di analisi secondo il cappello giallo.
\end{itemize}

% __________________________________________________%


% UC124 - Richiesta analisi secondo il cappello verde
\hypertarget{uc124}{}
\subsubsection{UC 124: Richiesta analisi secondo il cappello verde}

\begin{figure}[H]
    \centering
    \includegraphics[width=0.6\textwidth]{richiesta_cappello_verde.png}
    \caption{UC 124 – Richiesta analisi secondo il cappello verde}
\end{figure}

\begin{itemize}
    \item \textbf{Attore Principale}: Utente
    \item \textbf{Attore Secondario}: LLM
    \item \textbf{Precondizioni}: Una porzione di testo è selezionata nell'editor e il servizio AI è disponibile.
    \item \textbf{Scenario Principale}:
    \begin{enumerate}
        \item L'utente richiede un'analisi del testo selezionato secondo il cappello verde.
        \item Il sistema acquisisce la richiesta dell'utente.
    \end{enumerate}
    \item \textbf{Postcondizioni}: La richiesta di analisi secondo il cappello verde è stata acquisita dal sistema.
    \item \textbf{Trigger}: L'utente desidera esplorare nuove idee, alternative e prospettive creative del testo selezionato.
\end{itemize}


% UC125 - Generazione analisi secondo il cappello verde
\hypertarget{uc125}{}
\subsubsection{UC 125: Generazione analisi secondo il cappello verde}

\begin{figure}[H]
    \centering
    \includegraphics[width=0.6\textwidth]{generazione_cappello_verde.png}
    \caption{UC 125 – Generazione analisi secondo il cappello verde}
\end{figure}

\begin{itemize}
    \item \textbf{Attore Principale}: Utente
    \item \textbf{Attore Secondario}: LLM
    \item \textbf{Precondizioni}: La richiesta di analisi secondo il cappello verde è stata acquisita dal sistema.
    \item \textbf{Scenario Principale}:
    \begin{enumerate}
        \item Il sistema inoltra al servizio LLM la richiesta di analisi secondo il cappello verde.
        \item Il servizio LLM prende in carico la richiesta e avvia l'elaborazione.
    \end{enumerate}
    \item \textbf{Postcondizioni}: La richiesta di analisi secondo il cappello verde è in elaborazione presso il servizio LLM.
    \item \textbf{Trigger}: Il sistema ha acquisito la richiesta di analisi secondo il cappello verde.
\end{itemize}


% UC126 - Visualizzazione elaborazione analisi secondo il cappello verde
\hypertarget{uc126}{}
\subsubsection{UC 126: Visualizzazione elaborazione analisi secondo il cappello verde}

\begin{figure}[H]
    \centering
    \includegraphics[width=0.6\textwidth]{elaborazione_cappello_verde.png}
    \caption{UC 126 – Visualizzazione elaborazione analisi secondo il cappello verde}
\end{figure}

\begin{itemize}
    \item \textbf{Attore Principale}: Utente
    \item \textbf{Precondizioni}: La richiesta di analisi secondo il cappello verde è in elaborazione presso il servizio LLM.
    \item \textbf{Scenario Principale}:
    \begin{enumerate}
        \item Il sistema visualizza una finestra modale contenente un indicatore di caricamento.
        \item Il sistema rende disponibile il comando ``Interrompi''.
    \end{enumerate}
    \item \textbf{Postcondizioni}: La finestra modale di elaborazione è visualizzata e la richiesta è ancora in corso.
    \item \textbf{Trigger}: Il servizio LLM ha avviato l'elaborazione della richiesta di analisi secondo il cappello verde.
\end{itemize}


% UC127 - Interruzione elaborazione analisi secondo il cappello verde
\hypertarget{uc127}{}
\subsubsection{UC 127: Interruzione elaborazione analisi secondo il cappello verde}

\begin{figure}[H]
    \centering
    \includegraphics[width=0.6\textwidth]{interruzione_cappello_verde.png}
    \caption{UC 127 – Interruzione elaborazione analisi secondo il cappello verde}
\end{figure}

\begin{itemize}
    \item \textbf{Attore Principale}: Utente
    \item \textbf{Attore Secondario}: LLM
    \item \textbf{Precondizioni}: La finestra modale di elaborazione è visualizzata e la richiesta è ancora in corso.
    \item \textbf{Scenario Principale}:
    \begin{enumerate}
        \item L'utente attiva il comando ``Interrompi''.
        \item Il sistema richiede al servizio LLM l'interruzione dell'elaborazione.
        \item Il sistema chiude la finestra modale di elaborazione.
    \end{enumerate}
    \item \textbf{Postcondizioni}: L'elaborazione è interrotta e il contenuto della nota rimane invariato.
    \item \textbf{Trigger}: L'utente desidera interrompere la richiesta ancora in elaborazione.
\end{itemize}


% UC128 - Visualizzazione proposta analisi secondo il cappello verde
\hypertarget{uc128}{}
\subsubsection{UC 128: Visualizzazione proposta analisi secondo il cappello verde}

\begin{figure}[H]
    \centering
    \includegraphics[width=0.6\textwidth]{proposta_cappello_verde.png}
    \caption{UC 128 – Visualizzazione proposta analisi secondo il cappello verde}
\end{figure}

\begin{itemize}
    \item \textbf{Attore Principale}: Utente
    \item \textbf{Precondizioni}: Il servizio LLM ha completato la generazione della proposta di analisi secondo il cappello verde.
    \item \textbf{Scenario Principale}:
    \begin{enumerate}
        \item Il sistema chiude la finestra modale di elaborazione.
        \item Il sistema visualizza l'analisi secondo il cappello verde in una finestra modale dedicata.
        \item Il sistema rende disponibili i comandi ``Accetta'', ``Rifiuta'' e ``Rigenera''.
    \end{enumerate}
    \item \textbf{Postcondizioni}: La proposta è visualizzata e il contenuto della nota rimane invariato.
    \item \textbf{Trigger}: Il servizio LLM ha restituito la proposta di analisi secondo il cappello verde.
\end{itemize}


% UC129 - Accettazione proposta analisi secondo il cappello verde
\hypertarget{uc129}{}
\subsubsection{UC 129: Accettazione proposta analisi secondo il cappello verde}

\begin{figure}[H]
    \centering
    \includegraphics[width=0.6\textwidth]{accettazione_cappello_verde.png}
    \caption{UC 129 – Accettazione proposta analisi secondo il cappello verde}
\end{figure}

\begin{itemize}
    \item \textbf{Attore Principale}: Utente
    \item \textbf{Precondizioni}: La finestra modale contenente la proposta di analisi secondo il cappello verde è visualizzata.
    \item \textbf{Scenario Principale}:
    \begin{enumerate}
        \item L'utente attiva il comando ``Accetta''.
        \item Il sistema sostituisce il testo inizialmente selezionato con l'analisi proposta secondo il cappello verde.
        \item Il sistema chiude la finestra modale della proposta.
    \end{enumerate}
    \item \textbf{Postcondizioni}: Il testo inizialmente selezionato è sostituito con l'analisi secondo il cappello verde accettata e la finestra modale è chiusa.
    \item \textbf{Trigger}: L'utente desidera applicare la proposta visualizzata.
\end{itemize}


% UC130 - Rifiuto proposta analisi secondo il cappello verde
\hypertarget{uc130}{}
\subsubsection{UC 130: Rifiuto proposta analisi secondo il cappello verde}

\begin{figure}[H]
    \centering
    \includegraphics[width=0.6\textwidth]{rifiuto_cappello_verde.png}
    \caption{UC 130 – Rifiuto proposta analisi secondo il cappello verde}
\end{figure}

\begin{itemize}
    \item \textbf{Attore Principale}: Utente
    \item \textbf{Precondizioni}: La finestra modale contenente la proposta di analisi secondo il cappello verde è visualizzata.
    \item \textbf{Scenario Principale}:
    \begin{enumerate}
        \item L'utente attiva il comando ``Rifiuta''.
        \item Il sistema scarta la proposta visualizzata.
        \item Il sistema chiude la finestra modale della proposta.
    \end{enumerate}
    \item \textbf{Postcondizioni}: La proposta è scartata, la finestra modale è chiusa e il contenuto della nota rimane invariato.
    \item \textbf{Trigger}: L'utente non desidera applicare la proposta visualizzata.
\end{itemize}


% UC131 - Rigenerazione proposta analisi secondo il cappello verde
\hypertarget{uc131}{}
\subsubsection{UC 131: Rigenerazione proposta analisi secondo il cappello verde}

\begin{figure}[H]
    \centering
    \includegraphics[width=0.6\textwidth]{rigenerazione_cappello_verde.png}
    \caption{UC 131 – Rigenerazione proposta analisi secondo il cappello verde}
\end{figure}

\begin{itemize}
    \item \textbf{Attore Principale}: Utente
    \item \textbf{Attore Secondario}: LLM
    \item \textbf{Precondizioni}: La finestra modale contenente la proposta di analisi secondo il cappello verde è visualizzata.
    \item \textbf{Scenario Principale}:
    \begin{enumerate}
        \item L'utente attiva il comando ``Rigenera''.
        \item Il sistema scarta la proposta visualizzata.
        \item Il sistema inoltra al servizio LLM una nuova richiesta di analisi secondo il cappello verde.
        \item Il sistema visualizza nuovamente la finestra modale di elaborazione.
    \end{enumerate}
    \item \textbf{Postcondizioni}: Una nuova richiesta è in elaborazione e il contenuto della nota rimane invariato.
    \item \textbf{Trigger}: L'utente desidera ottenere una proposta differente.
\end{itemize}


% UC132 - Visualizzazione errore generazione analisi secondo il cappello verde
\hypertarget{uc132}{}
\subsubsection{UC 132: Visualizzazione errore generazione analisi secondo il cappello verde}

\begin{figure}[H]
    \centering
    \includegraphics[width=0.6\textwidth]{errore_cappello_verde.png}
    \caption{UC 132 – Visualizzazione errore generazione analisi secondo il cappello verde}
\end{figure}

\begin{itemize}
    \item \textbf{Attore Principale}: Utente
    \item \textbf{Precondizioni}: È in corso una richiesta di analisi secondo il cappello verde.
    \item \textbf{Scenario Principale}:
    \begin{enumerate}
        \item Si verifica un errore durante la comunicazione con il servizio LLM o durante l'elaborazione della richiesta.
        \item Il sistema chiude la finestra modale di elaborazione.
        \item Il sistema visualizza un messaggio che segnala il mancato completamento della richiesta.
    \end{enumerate}
    \item \textbf{Postcondizioni}: L'utente è informato dell'errore e il contenuto della nota rimane invariato.
    \item \textbf{Trigger}: Si verifica un errore prima del completamento della richiesta di analisi secondo il cappello verde.
\end{itemize}

% __________________________________________________%


% UC133 - Richiesta analisi secondo il cappello blu
\hypertarget{uc133}{}
\subsubsection{UC 133: Richiesta analisi secondo il cappello blu}

\begin{figure}[H]
    \centering
    \includegraphics[width=0.6\textwidth]{richiesta_cappello_blu.png}
    \caption{UC 133 – Richiesta analisi secondo il cappello blu}
\end{figure}

\begin{itemize}
    \item \textbf{Attore Principale}: Utente
    \item \textbf{Attore Secondario}: LLM
    \item \textbf{Precondizioni}: Una porzione di testo è selezionata nell'editor e il servizio AI è disponibile.
    \item \textbf{Scenario Principale}:
    \begin{enumerate}
        \item L'utente richiede un'analisi del testo selezionato secondo il cappello blu.
        \item Il sistema acquisisce la richiesta dell'utente.
    \end{enumerate}
    \item \textbf{Postcondizioni}: La richiesta di analisi secondo il cappello blu è stata acquisita dal sistema.
    \item \textbf{Trigger}: L'utente desidera analizzare la coerenza logica e la struttura del testo selezionato.
\end{itemize}


% UC134 - Generazione analisi secondo il cappello blu
\hypertarget{uc134}{}
\subsubsection{UC 134: Generazione analisi secondo il cappello blu}

\begin{figure}[H]
    \centering
    \includegraphics[width=0.6\textwidth]{generazione_cappello_blu.png}
    \caption{UC 134 – Generazione analisi secondo il cappello blu}
\end{figure}

\begin{itemize}
    \item \textbf{Attore Principale}: Utente
    \item \textbf{Attore Secondario}: LLM
    \item \textbf{Precondizioni}: La richiesta di analisi secondo il cappello blu è stata acquisita dal sistema.
    \item \textbf{Scenario Principale}:
    \begin{enumerate}
        \item Il sistema inoltra al servizio LLM la richiesta di analisi secondo il cappello blu.
        \item Il servizio LLM prende in carico la richiesta e avvia l'elaborazione.
    \end{enumerate}
    \item \textbf{Postcondizioni}: La richiesta di analisi secondo il cappello blu è in elaborazione presso il servizio LLM.
    \item \textbf{Trigger}: Il sistema ha acquisito la richiesta di analisi secondo il cappello blu.
\end{itemize}


% UC135 - Visualizzazione elaborazione analisi secondo il cappello blu
\hypertarget{uc135}{}
\subsubsection{UC 135: Visualizzazione elaborazione analisi secondo il cappello blu}

\begin{figure}[H]
    \centering
    \includegraphics[width=0.6\textwidth]{elaborazione_cappello_blu.png}
    \caption{UC 135 – Visualizzazione elaborazione analisi secondo il cappello blu}
\end{figure}

\begin{itemize}
    \item \textbf{Attore Principale}: Utente
    \item \textbf{Precondizioni}: La richiesta di analisi secondo il cappello blu è in elaborazione presso il servizio LLM.
    \item \textbf{Scenario Principale}:
    \begin{enumerate}
        \item Il sistema visualizza una finestra modale contenente un indicatore di caricamento.
        \item Il sistema rende disponibile il comando ``Interrompi''.
    \end{enumerate}
    \item \textbf{Postcondizioni}: La finestra modale di elaborazione è visualizzata e la richiesta è ancora in corso.
    \item \textbf{Trigger}: Il servizio LLM ha avviato l'elaborazione della richiesta di analisi secondo il cappello blu.
\end{itemize}


% UC136 - Interruzione elaborazione analisi secondo il cappello blu
\hypertarget{uc136}{}
\subsubsection{UC 136: Interruzione elaborazione analisi secondo il cappello blu}

\begin{figure}[H]
    \centering
    \includegraphics[width=0.6\textwidth]{interruzione_cappello_blu.png}
    \caption{UC 136 – Interruzione elaborazione analisi secondo il cappello blu}
\end{figure}

\begin{itemize}
    \item \textbf{Attore Principale}: Utente
    \item \textbf{Attore Secondario}: LLM
    \item \textbf{Precondizioni}: La finestra modale di elaborazione è visualizzata e la richiesta è ancora in corso.
    \item \textbf{Scenario Principale}:
    \begin{enumerate}
        \item L'utente attiva il comando ``Interrompi''.
        \item Il sistema richiede al servizio LLM l'interruzione dell'elaborazione.
        \item Il sistema chiude la finestra modale di elaborazione.
    \end{enumerate}
    \item \textbf{Postcondizioni}: L'elaborazione è interrotta e il contenuto della nota rimane invariato.
    \item \textbf{Trigger}: L'utente desidera interrompere la richiesta ancora in elaborazione.
\end{itemize}


% UC137 - Visualizzazione proposta analisi secondo il cappello blu
\hypertarget{uc137}{}
\subsubsection{UC 137: Visualizzazione proposta analisi secondo il cappello blu}

\begin{figure}[H]
    \centering
    \includegraphics[width=0.6\textwidth]{proposta_cappello_blu.png}
    \caption{UC 137 – Visualizzazione proposta analisi secondo il cappello blu}
\end{figure}

\begin{itemize}
    \item \textbf{Attore Principale}: Utente
    \item \textbf{Precondizioni}: Il servizio LLM ha completato la generazione della proposta di analisi secondo il cappello blu.
    \item \textbf{Scenario Principale}:
    \begin{enumerate}
        \item Il sistema chiude la finestra modale di elaborazione.
        \item Il sistema visualizza l'analisi secondo il cappello blu in una finestra modale dedicata.
        \item Il sistema rende disponibili i comandi ``Accetta'', ``Rifiuta'' e ``Rigenera''.
    \end{enumerate}
    \item \textbf{Postcondizioni}: La proposta è visualizzata e il contenuto della nota rimane invariato.
    \item \textbf{Trigger}: Il servizio LLM ha restituito la proposta di analisi secondo il cappello blu.
\end{itemize}


% UC138 - Accettazione proposta analisi secondo il cappello blu
\hypertarget{uc138}{}
\subsubsection{UC 138: Accettazione proposta analisi secondo il cappello blu}

\begin{figure}[H]
    \centering
    \includegraphics[width=0.6\textwidth]{accettazione_cappello_blu.png}
    \caption{UC 138 – Accettazione proposta analisi secondo il cappello blu}
\end{figure}

\begin{itemize}
    \item \textbf{Attore Principale}: Utente
    \item \textbf{Precondizioni}: La finestra modale contenente la proposta di analisi secondo il cappello blu è visualizzata.
    \item \textbf{Scenario Principale}:
    \begin{enumerate}
        \item L'utente attiva il comando ``Accetta''.
        \item Il sistema sostituisce il testo inizialmente selezionato con l'analisi proposta secondo il cappello blu.
        \item Il sistema chiude la finestra modale della proposta.
    \end{enumerate}
    \item \textbf{Postcondizioni}: Il testo inizialmente selezionato è sostituito con l'analisi secondo il cappello blu accettata e la finestra modale è chiusa.
    \item \textbf{Trigger}: L'utente desidera applicare la proposta visualizzata.
\end{itemize}


% UC139 - Rifiuto proposta analisi secondo il cappello blu
\hypertarget{uc139}{}
\subsubsection{UC 139: Rifiuto proposta analisi secondo il cappello blu}

\begin{figure}[H]
    \centering
    \includegraphics[width=0.6\textwidth]{rifiuto_cappello_blu.png}
    \caption{UC 139 – Rifiuto proposta analisi secondo il cappello blu}
\end{figure}

\begin{itemize}
    \item \textbf{Attore Principale}: Utente
    \item \textbf{Precondizioni}: La finestra modale contenente la proposta di analisi secondo il cappello blu è visualizzata.
    \item \textbf{Scenario Principale}:
    \begin{enumerate}
        \item L'utente attiva il comando ``Rifiuta''.
        \item Il sistema scarta la proposta visualizzata.
        \item Il sistema chiude la finestra modale della proposta.
    \end{enumerate}
    \item \textbf{Postcondizioni}: La proposta è scartata, la finestra modale è chiusa e il contenuto della nota rimane invariato.
    \item \textbf{Trigger}: L'utente non desidera applicare la proposta visualizzata.
\end{itemize}


% UC140 - Rigenerazione proposta analisi secondo il cappello blu
\hypertarget{uc140}{}
\subsubsection{UC 140: Rigenerazione proposta analisi secondo il cappello blu}

\begin{figure}[H]
    \centering
    \includegraphics[width=0.6\textwidth]{rigenerazione_cappello_blu.png}
    \caption{UC 140 – Rigenerazione proposta analisi secondo il cappello blu}
\end{figure}

\begin{itemize}
    \item \textbf{Attore Principale}: Utente
    \item \textbf{Attore Secondario}: LLM
    \item \textbf{Precondizioni}: La finestra modale contenente la proposta di analisi secondo il cappello blu è visualizzata.
    \item \textbf{Scenario Principale}:
    \begin{enumerate}
        \item L'utente attiva il comando ``Rigenera''.
        \item Il sistema scarta la proposta visualizzata.
        \item Il sistema inoltra al servizio LLM una nuova richiesta di analisi secondo il cappello blu.
        \item Il sistema visualizza nuovamente la finestra modale di elaborazione.
    \end{enumerate}
    \item \textbf{Postcondizioni}: Una nuova richiesta è in elaborazione e il contenuto della nota rimane invariato.
    \item \textbf{Trigger}: L'utente desidera ottenere una proposta differente.
\end{itemize}


% UC141 - Visualizzazione errore generazione analisi secondo il cappello blu
\hypertarget{uc141}{}
\subsubsection{UC 141: Visualizzazione errore generazione analisi secondo il cappello blu}

\begin{figure}[H]
    \centering
    \includegraphics[width=0.6\textwidth]{errore_cappello_blu.png}
    \caption{UC 141 – Visualizzazione errore generazione analisi secondo il cappello blu}
\end{figure}

\begin{itemize}
    \item \textbf{Attore Principale}: Utente
    \item \textbf{Precondizioni}: È in corso una richiesta di analisi secondo il cappello blu.
    \item \textbf{Scenario Principale}:
    \begin{enumerate}
        \item Si verifica un errore durante la comunicazione con il servizio LLM o durante l'elaborazione della richiesta.
        \item Il sistema chiude la finestra modale di elaborazione.
        \item Il sistema visualizza un messaggio che segnala il mancato completamento della richiesta.
    \end{enumerate}
    \item \textbf{Postcondizioni}: L'utente è informato dell'errore e il contenuto della nota rimane invariato.
    \item \textbf{Trigger}: Si verifica un errore prima del completamento della richiesta di analisi secondo il cappello blu.
\end{itemize}

% __________________________________________________%

% --------------------------------------------------------------------------
% 3. Macroarea 3: Gestione I/O e Persistenza
% --------------------------------------------------------------------------
\newpage
\subsection{Macroarea 3 — Gestione I/O e Persistenza}
\textit{Importazione, esportazione e persistenza delle note.}

\vspace{0.3cm}
\noindent
\fcolorbox{primarycolor!50}{primarycolor!5}{%
\begin{minipage}{\dimexpr\textwidth-2\fboxsep-2\fboxrule\relax}
\vspace{0.2cm}
\textbf{\textcolor{primarycolor}{Nota Esplicativa}} \\
\vspace{-0.1cm}

Una nota aperta nell'editor è modificabile tramite tutte le funzionalità descritte nella Macroarea 1 e nella Macroarea 2. 
Le funzionalità di salvataggio, importazione, esportazione e persistenza operano sul contenuto corrente della nota aperta nell'editor.
\vspace{0.1cm}
\end{minipage}%
}
\vspace{0.3cm}

% UC144 - Creazione di una nuova nota vuota
\hypertarget{uc144}{}
\subsubsection{UC 144: (Opzionale) Creazione di una nuova nota vuota}

% use case diagram
\begin{figure}[H]
    \centering
    \includegraphics[width=0.5\textwidth]{creazione_nota.png}
    \caption{UC 144 – Creazione di una nuova nota}
\end{figure}

\begin{itemize}
    \item \textbf{Attore Principale}: Utente
    \item \textbf{Precondizioni}: Il sistema è avviato e l'interfaccia è correttamente caricata.
    \item \textbf{Scenario Principale}:
    \begin{enumerate}
        \item L'utente attiva il comando ``Nuova Nota''.
        \item Il sistema richiede l'inserimento di un titolo per la nota.
        \item L'utente inserisce il titolo desiderato.
        \item Il sistema crea una nuova nota associando il titolo inserito.
        \item Il sistema apre la nota nell'editor.
    \end{enumerate}
    \item \textbf{Postcondizioni}: Una nuova nota con titolo associato è aperta e modificabile nell'editor.
    \item \textbf{Trigger}: L'utente desidera creare una nuova nota.
\end{itemize}

% UC145 - Eliminazione di una nota
\hypertarget{uc145}{}
\subsubsection{UC 145: (Opzionale) Eliminazione di una nota}

% use case diagram
\begin{figure}[H]
    \centering
    \includegraphics[width=0.5\textwidth]{eliminazione_nota.png}
    \caption{UC 145 – Eliminazione di una nota}
\end{figure}

\begin{itemize}
    \item \textbf{Attore Principale}: Utente
    \item \textbf{Precondizioni}: Il sistema è avviato ed è presente almeno una nota.
    \item \textbf{Scenario Principale}:
    \begin{enumerate}
        \item L'utente seleziona la nota da eliminare.
        \item Il sistema richiede conferma.
        \item L'utente conferma.
        \item Il sistema elimina la nota.
    \end{enumerate}
    \item \textbf{Postcondizioni}: La nota selezionata è rimossa dal sistema.
    \item \textbf{Trigger}: L'utente decide di rimuovere una nota non più necessaria.
    
    \item \textbf{Inclusioni}:
    \begin{itemize}
        \item \hyperlink{uc145.1}{UC 145.1: Conferma eliminazione}
    \end{itemize}

    \item \textbf{Estensioni}: 
    \begin{itemize}
        \item \hyperlink{uc145.2}{UC 145.2: Annullamento eliminazione} $\rightarrow$ Condizione: risposta negativa al dialogo di conferma.
        \item \hyperlink{uc145.3}{UC 145.3: Visualizzazione errore eliminazione} $\rightarrow$ Condizione: errore tecnico durante l'esecuzione dell'eliminazione al passo 4.
    \end{itemize}
\end{itemize}

% UC145.1 - Conferma eliminazione
\hypertarget{uc145.1}{}
\subsubsection{UC 145.1: Conferma eliminazione}

% use case diagram
\begin{figure}[H]
    \centering
    \includegraphics[width=0.7\textwidth]{conferma_eliminazione.png}
    \caption{UC 145.1 – Conferma eliminazione}
\end{figure}

\begin{itemize}
    \item \textbf{Attore Principale}: Utente
    \item \textbf{Precondizioni}: L'utente ha richiesto l'eliminazione di una nota.
    \item \textbf{Scenario Principale}:
    \begin{enumerate}
        \item Il sistema mostra un dialogo di conferma.
        \item L'utente conferma l'operazione.
    \end{enumerate}
    \item \textbf{Postcondizioni}: La richiesta di eliminazione è confermata e il sistema procede.
    \item \textbf{Trigger}: Il sistema richiede una conferma esplicita prima di un'operazione distruttiva.
\end{itemize}

% UC145.2 - Annullamento eliminazione
\hypertarget{uc145.2}{}
\subsubsection{UC 145.2: Annullamento eliminazione}

% use case diagram
\begin{figure}[H]
    \centering
    \includegraphics[width=0.7\textwidth]{annullamento_eliminazione.png}
    \caption{UC 145.2 – Annullamento eliminazione}
\end{figure}

\begin{itemize}
    \item \textbf{Attore Principale}: Utente
    \item \textbf{Precondizioni}: Il dialogo di conferma eliminazione è visualizzato a schermo.
    \item \textbf{Scenario Principale}:
    \begin{enumerate}
        \item L'utente preme ``Annulla'' nel dialogo di conferma.
        \item Il sistema chiude il dialogo senza apportare modifiche.
    \end{enumerate}
    \item \textbf{Postcondizioni}: La nota rimane invariata; nessuna modifica è avvenuta.
    \item \textbf{Trigger}: L'utente decide di non procedere con l'eliminazione.
\end{itemize}

% UC145.3 - Visualizzazione errore eliminazione
\hypertarget{uc145.3}{}
\subsubsection{UC 145.3: Visualizzazione errore eliminazione}

% use case diagram
\begin{figure}[H]
    \centering
    \includegraphics[width=0.6\textwidth]{errore_eliminazione.png}
    \caption{UC 145.3 – Visualizzazione errore eliminazione}
\end{figure}

\begin{itemize}
    \item \textbf{Attore Principale}: Utente
    \item \textbf{Precondizioni}: L'utente ha confermato l'eliminazione della nota.
    \item \textbf{Scenario Principale}:
    \begin{enumerate}
        \item Il sistema riscontra un problema tecnico durante la rimozione della nota.
        \item Il sistema mostra un messaggio di errore all'utente.
    \end{enumerate}
    \item \textbf{Postcondizioni}: La nota non viene eliminata e l'utente viene informato dell'errore.
    \item \textbf{Trigger}: Il sistema operativo o il file system\textsubscript{G} impediscono la cancellazione del file.
\end{itemize}

% UC146 - Salvataggio della nota in formato Markdown
\hypertarget{uc146}{}
\subsubsection{UC 146: Salvataggio della nota in formato Markdown}

% use case diagram
\begin{figure}[H]
    \centering
    \includegraphics[width=0.5\textwidth]{download_nota.png}
    \caption{UC 146 – Salvataggio della nota in formato Markdown}
\end{figure}

\begin{itemize}
    \item \textbf{Attore Principale}: Utente
    \item \textbf{Precondizioni}: Una nota è aperta nell'editor.
    \item \textbf{Scenario Principale}:
    \begin{enumerate}
        \item L'utente seleziona la funzione di salvataggio della nota.
        \item Il sistema genera un file Markdown contenente il contenuto della nota corrente.
        \item Il sistema avvia il salvataggio del file sul dispositivo dell'utente.
    \end{enumerate}
    \item \textbf{Postcondizioni}: Una copia locale della nota viene salvata sul dispositivo dell'utente.
    \item \textbf{Trigger}: L'utente desidera salvare localmente la nota corrente.
    \item \textbf{Estensioni}:
    \begin{itemize}
        \item \hyperlink{uc146.1}{UC 146.1: Visualizzazione errore salvataggio} $ \rightarrow $ Condizione: si verifica un errore durante la procedura di salvataggio.
    \end{itemize}
\end{itemize}

% UC146.1 - Visualizzazione errore salvataggio
\hypertarget{uc146.1}{}
\subsubsection{UC 146.1: Visualizzazione errore salvataggio}\

\begin{figure}[H]
    \centering
    \includegraphics[width=0.6\textwidth]{errore_download.png}
    \caption{UC 146.1 – Visualizzazione errore salvataggio}
\end{figure}

\begin{itemize}
    \item \textbf{Attore Principale}: Utente
    \item \textbf{Precondizioni}: È in corso il salvataggio della nota.
    \item \textbf{Scenario Principale}:
    \begin{enumerate}
        \item Il sistema rileva un errore durante la generazione o il salvataggio del file.
        \item Il sistema mostra un messaggio di errore all'utente.
    \end{enumerate}
    \item \textbf{Postcondizioni}: Il file non viene salvato.
    \item \textbf{Trigger}: Errore durante la procedura di salvataggio.
\end{itemize}

% UC147 - Importazione nota locale
\hypertarget{uc147}{}
\subsubsection{UC 147: Importazione di una nota locale}

% use case diagram
\begin{figure}[H]
    \centering
    \includegraphics[width=0.5\textwidth]{upload_nota.png}
    \caption{UC 147 – Importazione di una nota locale}
\end{figure}

\begin{itemize}
    \item \textbf{Attore Principale}: Utente
    \item \textbf{Precondizioni}: L'utente possiede un file Markdown compatibile.
    \item \textbf{Scenario Principale}:
    \begin{enumerate}
        \item L'utente attiva la funzione di importazione nota.
        \item L'utente seleziona un file Markdown dal proprio dispositivo.
        \item Il sistema verifica la compatibilità del file selezionato.
        \item Il sistema importa il contenuto della nota nell'editor.
    \end{enumerate}
    \item \textbf{Postcondizioni}: La nota importata è aperta nell'editor.
    \item \textbf{Trigger}: L'utente desidera aprire una nota salvata localmente.
    \item \textbf{Estensioni}:
    \begin{itemize}
        \item \hyperlink{uc147.1}{UC 147.1: Visualizzazione errore formato file} $ \rightarrow $ Condizione: il file selezionato non è compatibile o è corrotto.
    \end{itemize}
\end{itemize}

% UC147.1 - Visualizzazione errore formato file
\hypertarget{uc147.1}{}
\subsubsection{UC 147.1: Visualizzazione errore formato file}

% use case diagram
\begin{figure}[H]
    \centering
    \includegraphics[width=0.6\textwidth]{errore_formato_file.png}
    \caption{UC 147.1 – Visualizzazione errore formato file}
\end{figure}

\begin{itemize}
    \item \textbf{Attore Principale}: Utente
    \item \textbf{Precondizioni}: È in corso una procedura di importazione.
    \item \textbf{Scenario Principale}:
    \begin{enumerate}
        \item Il sistema rileva un file incompatibile o corrotto.
        \item Il sistema mostra un messaggio di errore.
    \end{enumerate}
    \item \textbf{Postcondizioni}: La nota non viene importata.
    \item \textbf{Trigger}: L'utente seleziona un file non valido.
\end{itemize}

% UC148 - Esportazione nota in formato PDF
\hypertarget{uc148}{}
\subsubsection{UC 148: Esportazione nota in formato PDF}

\begin{figure}[H]
    \centering
    \includegraphics[width=0.5\textwidth]{esportazione_pdf.png}
    \caption{UC 148 – Esportazione nota in formato PDF}
\end{figure}

\begin{itemize}
    \item \textbf{Attore Principale}: Utente
    \item \textbf{Precondizioni}: Una nota è aperta nell'editor.
    \item \textbf{Scenario Principale}:
    \begin{enumerate}
        \item L'utente seleziona la funzione di esportazione in formato PDF.
        \item Il sistema converte la nota in formato PDF.
        \item Il sistema avvia il download del file esportato.
    \end{enumerate}
    \item \textbf{Postcondizioni}: L'utente ottiene una copia della nota in formato PDF.
    \item \textbf{Trigger}: L'utente desidera esportare la nota corrente come documento PDF.
    \item \textbf{Estensioni}:
    \begin{itemize}
        \item \hyperlink{uc148.1}{UC 148.1: Visualizzazione errore esportazione PDF} $ \rightarrow $ Condizione: si verifica un errore durante la conversione o il salvataggio in PDF.
    \end{itemize}
\end{itemize}

% UC148.1 - Visualizzazione errore esportazione PDF
\hypertarget{uc148.1}{}
\subsubsection{UC 148.1: Visualizzazione errore esportazione PDF}

\begin{figure}[H]
    \centering
    \includegraphics[width=0.6\textwidth]{errore_esportazione_pdf.png}
    \caption{UC 148.1 – Visualizzazione errore esportazione PDF}
\end{figure}

\begin{itemize}
    \item \textbf{Attore Principale}: Utente
    \item \textbf{Precondizioni}: È in corso una procedura di esportazione in PDF.
    \item \textbf{Scenario Principale}:
    \begin{enumerate}
        \item Il sistema rileva un errore durante la conversione o il download del file PDF.
        \item Il sistema mostra un messaggio di errore all'utente.
    \end{enumerate}
    \item \textbf{Postcondizioni}: L'esportazione in PDF non viene completata.
    \item \textbf{Trigger}: Errore durante la conversione del file in PDF.
\end{itemize}

% UC149 - Esportazione nota in formato HTML
\hypertarget{uc149}{}
\subsubsection{UC 149: Esportazione nota in formato HTML}

\begin{figure}[H]
    \centering
    \includegraphics[width=0.5\textwidth]{esportazione_html.png}
    \caption{UC 149 – Esportazione nota in formato HTML}
\end{figure}

\begin{itemize}
    \item \textbf{Attore Principale}: Utente
    \item \textbf{Precondizioni}: Una nota è aperta nell'editor.
    \item \textbf{Scenario Principale}:
    \begin{enumerate}
        \item L'utente seleziona la funzione di esportazione in formato HTML.
        \item Il sistema converte la nota in formato HTML.
        \item Il sistema avvia il download del file esportato.
    \end{enumerate}
    \item \textbf{Postcondizioni}: L'utente ottiene una copia della nota in formato HTML.
    \item \textbf{Trigger}: L'utente desidera esportare la nota corrente come pagina web HTML.
    \item \textbf{Estensioni}:
    \begin{itemize}
        \item \hyperlink{uc149.1}{UC 149.1: Visualizzazione errore esportazione HTML} $ \rightarrow $ Condizione: si verifica un errore durante la conversione o il salvataggio in HTML.
    \end{itemize}
\end{itemize}

% UC149.1 - Visualizzazione errore esportazione HTML
\hypertarget{uc149.1}{}
\subsubsection{UC 149.1: Visualizzazione errore esportazione HTML}

\begin{figure}[H]
    \centering
    \includegraphics[width=0.6\textwidth]{errore_esportazione_html.png}
    \caption{UC 149.1 – Visualizzazione errore esportazione HTML}
\end{figure}

\begin{itemize}
    \item \textbf{Attore Principale}: Utente
    \item \textbf{Precondizioni}: È in corso una procedura di esportazione in HTML.
    \item \textbf{Scenario Principale}:
    \begin{enumerate}
        \item Il sistema rileva un errore durante la conversione o il download del file HTML.
        \item Il sistema mostra un messaggio di errore all'utente.
    \end{enumerate}
    \item \textbf{Postcondizioni}: L'esportazione in HTML non viene completata.
    \item \textbf{Trigger}: Errore durante la conversione del file in HTML.
\end{itemize}

% UC150 - Esportazione nota in formato JSON
\hypertarget{uc150}{}
\subsubsection{UC 150: Esportazione nota in formato JSON}

\begin{figure}[H]
    \centering
    \includegraphics[width=0.5\textwidth]{esportazione_json.png}
    \caption{UC 150 – Esportazione nota in formato JSON}
\end{figure}

\begin{itemize}
    \item \textbf{Attore Principale}: Utente
    \item \textbf{Precondizioni}: Una nota è aperta nell'editor.
    \item \textbf{Scenario Principale}:
    \begin{enumerate}
        \item L'utente seleziona la funzione di esportazione in formato JSON.
        \item Il sistema converte i metadati e il contenuto della nota in formato JSON.
        \item Il sistema avvia il download del file esportato.
    \end{enumerate}
    \item \textbf{Postcondizioni}: L'utente ottiene una copia della nota in formato JSON.
    \item \textbf{Trigger}: L'utente desidera esportare la nota corrente come oggetto JSON.
    \item \textbf{Estensioni}:
    \begin{itemize}
        \item \hyperlink{uc150.1}{UC 150.1: Visualizzazione errore esportazione JSON} $ \rightarrow $ Condizione: si verifica un errore durante la conversione o il salvataggio in JSON.
    \end{itemize}
\end{itemize}

% UC150.1 - Visualizzazione errore esportazione JSON
\hypertarget{uc150.1}{}
\subsubsection{UC 150.1: Visualizzazione errore esportazione JSON}

\begin{figure}[H]
    \centering
    \includegraphics[width=0.6\textwidth]{errore_esportazione_json.png}
    \caption{UC 150.1 – Visualizzazione errore esportazione JSON}
\end{figure}

\begin{itemize}
    \item \textbf{Attore Principale}: Utente
    \item \textbf{Precondizioni}: È in corso una procedura di esportazione in JSON.
    \item \textbf{Scenario Principale}:
    \begin{enumerate}
        \item Il sistema rileva un errore durante la generazione o il download del file JSON.
        \item Il sistema mostra un messaggio di errore all'utente.
    \end{enumerate}
    \item \textbf{Postcondizioni}: L'esportazione in JSON non viene completata.
    \item \textbf{Trigger}: Errore durante la conversione del file in JSON.
\end{itemize}

% --------------------------------------------------------------------------
% Funzionalità opzionali di persistenza remota
% --------------------------------------------------------------------------

% UC151 - Registrazione utente 
\hypertarget{uc151}{}
\subsubsection{UC 151: (Opzionale) Registrazione utente}

\begin{figure}[H]
    \centering
    \includegraphics[width=0.5\textwidth]{registrazione_utente.png}
    \caption{UC 151 – Registrazione utente}
\end{figure}

\begin{itemize}
    \item \textbf{Attore Principale}: Utente
    \item \textbf{Precondizioni}: Il sistema è avviato, il servizio di persistenza remota è disponibile e l’utente non è autenticato nel sistema.
    \item \textbf{Scenario Principale}:
    \begin{enumerate}
        \item L'utente seleziona la funzione di registrazione.
        \item Il sistema mostra il modulo di registrazione.
        \item L'utente inserisce le informazioni richieste (es.\ email e password).
        \item Il sistema valida i dati inseriti.
        \item Il sistema crea il nuovo account utente.
        \item Il sistema conferma l'avvenuta registrazione.
    \end{enumerate}
    \item \textbf{Postcondizioni}: L'account utente viene creato correttamente.
    \item \textbf{Trigger}: L'utente desidera accedere alle funzionalità di sincronizzazione remota.
    \item \textbf{Estensioni}:
    \begin{itemize}
        \item \hyperlink{uc151.1}{UC 151.1: Visualizzazione errore email registrazione} $ \rightarrow $ Condizione: email non valida o già registrata al passo 4.
        \item \hyperlink{uc151.2}{UC 151.2: Visualizzazione errore password registrazione} $ \rightarrow $ Condizione: password non conforme ai requisiti al passo 4.
    \end{itemize}
\end{itemize}

% UC151.1 - Visualizzazione errore email registrazione
\hypertarget{uc151.1}{}
\subsubsection{UC 151.1: Visualizzazione errore email registrazione}\

\begin{figure}[H]
    \centering
    \includegraphics[width=0.6\textwidth]{errore_email_registrazione.png}
    \caption{UC 151.1 – Visualizzazione errore email registrazione}
\end{figure}

\begin{itemize}
    \item \textbf{Attore Principale}: Utente
    \item \textbf{Precondizioni}: È in corso la procedura di registrazione.
    \item \textbf{Scenario Principale}:
    \begin{enumerate}
        \item Il sistema rileva che l'email inserita non è valida oppure è già associata a un account esistente.
        \item Il sistema mostra un messaggio di errore all'utente.
    \end{enumerate}
    \item \textbf{Postcondizioni}: La registrazione non viene completata.
    \item \textbf{Trigger}: Inserimento di email non valida o già presente nel sistema.
\end{itemize}

% UC151.2 - Visualizzazione errore password registrazione
\hypertarget{uc151.2}{}
\subsubsection{UC 151.2: Visualizzazione errore password registrazione}\

\begin{figure}[H]
    \centering
    \includegraphics[width=0.6\textwidth]{errore_password_registrazione.png}
    \caption{UC 151.2 – Visualizzazione errore password registrazione}
\end{figure}

\begin{itemize}
    \item \textbf{Attore Principale}: Utente
    \item \textbf{Precondizioni}: È in corso la procedura di registrazione.
    \item \textbf{Scenario Principale}:
    \begin{enumerate}
        \item Il sistema rileva che la password inserita non soddisfa i requisiti minimi previsti.
        \item Il sistema mostra un messaggio di errore all'utente.
    \end{enumerate}
    \item \textbf{Postcondizioni}: La registrazione non viene completata.
    \item \textbf{Trigger}: Inserimento di password non conforme ai requisiti di sicurezza.
\end{itemize}

% UC152 - Autenticazione utente
\hypertarget{uc152}{}
\subsubsection{UC 152: Autenticazione utente}

\begin{figure}[H]
    \centering
    \includegraphics[width=0.5\textwidth]{autenticazione_utente.png}
    \caption{UC 152 – Autenticazione utente}
\end{figure}

\begin{itemize}
    \item \textbf{Attore Principale}: Utente
    \item \textbf{Precondizioni}: L’utente non è autenticato nel sistema e possiede un account registrato.
    \item \textbf{Scenario Principale}:
    \begin{enumerate}
        \item L'utente seleziona la funzione di accesso.
        \item Il sistema mostra il modulo di autenticazione.
        \item L'utente inserisce l'indirizzo email associato all'account.
        \item L'utente inserisce la password associata all'account.
        \item Il sistema verifica la corrispondenza tra indirizzo email e password.
        \item Il sistema autentica l'utente e abilita le funzionalità di persistenza remota.
    \end{enumerate}
    \item \textbf{Postcondizioni}: L'utente risulta autenticato nel sistema.
    \item \textbf{Trigger}: L'utente desidera accedere alle funzionalità remote.
    \item \textbf{Estensioni}:
    \begin{itemize}
        \item \hyperlink{uc152.1}{UC 152.1: Visualizzazione errore autenticazione} $ \rightarrow $ Condizione: credenziali non valide al passo 4.
        %\item \hyperlink{uc152.2}{UC 152.2: Visualizzazione errore utente non registrato} $ \rightarrow $ Condizione: utente non registrato al passo 4.
    \end{itemize}
\end{itemize}

% UC152.1 - Visualizzazione errore autenticazione
\hypertarget{uc152.1}{}
\subsubsection{UC 152.1: Visualizzazione errore autenticazione}\

\begin{figure}[H]
    \centering
    \includegraphics[width=0.6\textwidth]{errore_autenticazione.png}
    \caption{UC 152.1 – Visualizzazione errore autenticazione}
\end{figure}

\begin{itemize}
    \item \textbf{Attore Principale}: Utente
    \item \textbf{Precondizioni}: È in corso l'autenticazione dell'utente.
    \item \textbf{Scenario Principale}:
    \begin{enumerate}
        \item Il sistema rileva che le credenziali inserite non sono corrette.
        \item Il sistema mostra un messaggio di errore all'utente.
    \end{enumerate}
    \item \textbf{Postcondizioni}: L'utente non viene autenticato.
    \item \textbf{Trigger}: Inserimento di credenziali non valide.
\end{itemize}

% UC152.2 - Visualizzazione errore utente non registrato

% Utente non registrato è un tipo di credenziale non valida
\begin{comment}
    
\hypertarget{uc152.2}{}
\subsubsection{UC 152.2: Visualizzazione errore utente non registrato}\

\begin{figure}[H]
    \centering
    %\includegraphics[width=0.6\textwidth]{errore_utente_non_registrato.png}
    \caption{UC 152.2 – Visualizzazione errore utente non registrato}
\end{figure}

\begin{itemize}
    \item \textbf{Attore Principale}: Utente
    \item \textbf{Precondizioni}: È in corso l'autenticazione dell'utente.
    \item \textbf{Scenario Principale}:
    \begin{enumerate}
        \item Il sistema verifica che l'utente non possiede un account registrato.
        \item Il sistema informa l'utente che è necessaria la registrazione.
    \end{enumerate}
    \item \textbf{Postcondizioni}: L'accesso non viene consentito.
    \item \textbf{Trigger}: Tentativo di autenticazione con account inesistente.
\end{itemize}
\end{comment}

% UC153 - Logout utente
\hypertarget{uc153}{}
\subsubsection{UC 153: (Opzionale) Logout\textsubscript{G} utente}

\begin{figure}[H]
    \centering
    \includegraphics[width=0.5\textwidth]{logout_utente.png}
    \caption{UC 153 – Logout utente}
\end{figure}

\begin{itemize}
    \item \textbf{Attore Principale}: Utente
    \item \textbf{Precondizioni}: L'utente è autenticato nel sistema.
    \item \textbf{Scenario Principale}:
    \begin{enumerate}
        \item L'utente seleziona la funzione di logout.
        \item Il sistema termina la sessione utente autenticata.
        \item Il sistema disabilita le funzionalità di persistenza remota.
        \item Il sistema riporta l'utente alla modalità non autenticata.
    \end{enumerate}
    \item \textbf{Postcondizioni}: L'utente non risulta più autenticato nel sistema.
    \item \textbf{Trigger}: L'utente desidera terminare la sessione autenticata.
\end{itemize}

% UC154 - Visualizzazione elenco note archiviate
\hypertarget{uc154}{}
\subsubsection{UC 154: (Opzionale) Visualizzazione elenco note archiviate}

\begin{figure}[H]
    \centering
    \includegraphics[width=0.5\textwidth]{lista_note_salvate.png}
    \caption{UC 154 – Visualizzazione elenco note archiviate}
\end{figure}

\begin{itemize}
    \item \textbf{Attore Principale}: Utente
    \item \textbf{Precondizioni}: L'utente è autenticato e il database remoto è disponibile.
    \item \textbf{Scenario Principale}:
    \begin{enumerate}
        \item L'utente accede alla sezione delle note sincronizzate.
        \item Il sistema recupera le note associate all'account dell'utente dal database remoto.
        \item Il sistema mostra l'elenco delle note archiviate.
    \end{enumerate}
    \item \textbf{Postcondizioni}: L'utente visualizza la lista delle proprie note archiviate.
    \item \textbf{Trigger}: L'utente desidera consultare le note salvate nel database remoto.
    
    \item \textbf{Inclusioni}:
    \begin{itemize}
        \item \hyperlink{uc154.1}{UC 154.1: Visualizzazione titolo nota remota}
        \item \hyperlink{uc154.2}{UC 154.2: Visualizzazione data ultima modifica nota remota}
        \item \hyperlink{uc154.3}{UC 154.3: Visualizzazione anteprima contenuto nota remota}
    \end{itemize}

    \item \textbf{Estensioni}:
    \begin{itemize}
        \item \hyperlink{uc154.4}{UC 154.4: Visualizzazione errore recupero lista note remote} $ \rightarrow $ Condizione: errore durante il recupero delle note al passo 2.
    \end{itemize}
\end{itemize}

% UC154.1 - Visualizzazione titolo nota remota
\hypertarget{uc154.1}{}
\subsubsection{UC 154.1: Visualizzazione titolo nota remota}

\begin{figure}[H]
    \centering
    \includegraphics[width=0.6\textwidth]{visualizzazione_titolo_remota.png}
    \caption{UC 154.1 – Visualizzazione titolo nota remota}
\end{figure}

\begin{itemize}
    \item \textbf{Attore Principale}: Utente
    \item \textbf{Precondizioni}: Il sistema sta mostrando un elemento all'interno dell'elenco delle note archiviate.
    \item \textbf{Scenario Principale}:
    \begin{enumerate}
        \item Il sistema presenta il titolo testuale della nota recuperata dal database remoto.
    \end{enumerate}
    \item \textbf{Postcondizioni}: Il titolo della nota è visibile all'utente.
    \item \textbf{Trigger}: Generazione della lista delle note archiviate.
\end{itemize}

% UC154.2 - Visualizzazione data ultima modifica nota remota
\hypertarget{uc154.2}{}
\subsubsection{UC 154.2: Visualizzazione data ultima modifica nota remota}

\begin{figure}[H]
    \centering
    \includegraphics[width=0.6\textwidth]{visualizzazione_data_remota.png}
    \caption{UC 154.2 – Visualizzazione data ultima modifica nota remota}
\end{figure}

\begin{itemize}
    \item \textbf{Attore Principale}: Utente
    \item \textbf{Precondizioni}: Il sistema sta mostrando un elemento all'interno dell'elenco delle note archiviate.
    \item \textbf{Scenario Principale}:
    \begin{enumerate}
        \item Il sistema presenta il timestamp o la data formattata relativa all'ultimo salvataggio della nota nel database remoto.
    \end{enumerate}
    \item \textbf{Postcondizioni}: La data di ultima modifica della nota è visibile all'utente.
    \item \textbf{Trigger}: Generazione della lista delle note archiviate.
\end{itemize}

% UC154.3 - Visualizzazione anteprima contenuto nota remota
\hypertarget{uc154.3}{}
\subsubsection{UC 154.3: Visualizzazione anteprima contenuto nota remota}

\begin{figure}[H]
    \centering
    \includegraphics[width=0.6\textwidth]{visualizzazione_anteprima_remota.png}
    \caption{UC 154.3 – Visualizzazione anteprima contenuto nota remota}
\end{figure}

\begin{itemize}
    \item \textbf{Attore Principale}: Utente
    \item \textbf{Precondizioni}: Il sistema sta mostrando un elemento all'interno dell'elenco delle note archiviate.
    \item \textbf{Scenario Principale}:
    \begin{enumerate}
        \item Il sistema estrapola e presenta i primi caratteri del contenuto testuale della nota remota per fornire un breve riassunto visivo.
    \end{enumerate}
    \item \textbf{Postcondizioni}: Una porzione di anteprima del testo della nota è visibile all'utente.
    \item \textbf{Trigger}: Generazione della lista delle note archiviate.
\end{itemize}

% UC154.4 - Visualizzazione errore recupero lista note remote
\hypertarget{uc154.4}{}
\subsubsection{UC 154.4: Visualizzazione errore recupero lista note remote}
\begin{figure}[H]
    \centering
    \includegraphics[width=0.6\textwidth]{errore_lista_note_remote.png}
    \caption{UC 154.4 – Visualizzazione errore recupero lista note remote}
\end{figure}

\begin{itemize}
    \item \textbf{Attore Principale}: Utente
    \item \textbf{Precondizioni}: È in corso il recupero delle note remote.
    \item \textbf{Scenario Principale}:
    \begin{enumerate}
        \item Il sistema rileva un errore di comunicazione con il database PostgreSQL o un errore applicativo durante il recupero.
        \item Il sistema mostra un messaggio di errore all'utente.
    \end{enumerate}
    \item \textbf{Postcondizioni}: La lista delle note non viene visualizzata.
    \item \textbf{Trigger}: Problemi di connessione o indisponibilità del database remoto.
\end{itemize}

% UC155 - Salvataggio nota su database remoto
\hypertarget{uc155}{}
\subsubsection{UC 155: (Opzionale) Salvataggio nota su database remoto}

% use case diagram
\begin{figure}[H]
    \centering
    \includegraphics[width=0.5\textwidth]{salvataggio_db.png}
    \caption{UC 155 – Salvataggio nota su database remoto}
\end{figure}

\begin{itemize}
    \item \textbf{Attore Principale}: Utente
    \item \textbf{Precondizioni}: L'utente è autenticato e una nota è aperta nell'editor. Il sistema utilizza un database PostgreSQL configurato a livello applicativo per la memorizzazione remota delle note.
    \item \textbf{Scenario Principale}:
    \begin{enumerate}
        \item L'utente seleziona la funzione di salvataggio remoto.
        \item Il sistema sincronizza il contenuto della nota con il database remoto associato all'utente.
        \item Il sistema conferma l'avvenuto salvataggio.
    \end{enumerate}
    \item \textbf{Postcondizioni}: La nota viene salvata nel database remoto dell'utente.
    \item \textbf{Trigger}: L'utente desidera sincronizzare la nota corrente.
    \item \textbf{Estensioni}:
    \begin{itemize}
        \item \hyperlink{uc155.1}{UC 155.1: Visualizzazione errore sincronizzazione remota} $ \rightarrow $ Condizione: errore di sincronizzazione al passo 2.
    \end{itemize}
\end{itemize}

% UC155.1 - Visualizzazione errore sincronizzazione remota
\hypertarget{uc155.1}{}
\subsubsection{UC 155.1: Visualizzazione errore sincronizzazione remota}

% use case diagram
\begin{figure}[H]
    \centering
    \includegraphics[width=0.6\textwidth]{errore_salvataggio_db.png}
    \caption{UC 155.1 – Visualizzazione errore sincronizzazione remota}
\end{figure}

\begin{itemize}
    \item \textbf{Attore Principale}: Utente
    \item \textbf{Precondizioni}: È in corso il salvataggio remoto della nota.
    \item \textbf{Scenario Principale}:
    \begin{enumerate}
        \item Il sistema rileva un errore durante la comunicazione con il database PostgreSQL per la sincronizzazione della nota.
        \item Il sistema mostra un messaggio di errore all'utente.
    \end{enumerate}
    \item \textbf{Postcondizioni}: La nota non viene salvata nel database remoto.
    \item \textbf{Trigger}: Problemi di connessione o indisponibilità del database remoto.
\end{itemize}

% UC156 - Apertura nota da lista note salvate
\hypertarget{uc156}{}
\subsubsection{UC 156: (Opzionale) Apertura nota da lista note salvate}

\begin{figure}[H]
    \centering
    \includegraphics[width=0.5\textwidth]{apertura_nota_remota.png}
    \caption{UC 156 – Apertura nota da lista note salvate}
\end{figure}

\begin{itemize}
    \item \textbf{Attore Principale}: Utente
    \item \textbf{Precondizioni}: L'utente è autenticato e la lista delle note salvate è disponibile.
    \item \textbf{Scenario Principale}:
    \begin{enumerate}
        \item L'utente seleziona una nota dalla lista delle note salvate.
        \item Il sistema recupera il contenuto della nota selezionata.
        \item Il sistema carica la nota nell'editor.
    \end{enumerate}
    \item \textbf{Postcondizioni}: La nota selezionata è aperta nell'editor.
    \item \textbf{Trigger}: L'utente desidera aprire una nota salvata nel database remoto.
    \item \textbf{Estensioni}:
    \begin{itemize}
        \item \hyperlink{uc156.1}{UC 156.1: Visualizzazione errore apertura nota remota} $ \rightarrow $ Condizione: errore durante il recupero della nota al passo 2.
    \end{itemize}
\end{itemize}

% UC156.1 - Visualizzazione errore apertura nota remota
\hypertarget{uc156.1}{}
\subsubsection{UC 156.1: Visualizzazione errore apertura nota remota}\

\begin{figure}[H]
    \centering
    \includegraphics[width=0.6\textwidth]{errore_apertura_remota.png}
    \caption{UC 156.1 – Visualizzazione errore apertura nota remota}
\end{figure}

\begin{itemize}
    \item \textbf{Attore Principale}: Utente
    \item \textbf{Precondizioni}: È in corso il recupero della nota remota.
    \item \textbf{Scenario Principale}:
    \begin{enumerate}
        \item Il sistema non riesce a recuperare la nota selezionata.
        \item Il sistema mostra un messaggio di errore all'utente.
    \end{enumerate}
    \item \textbf{Postcondizioni}: L'editor rimane invariato.
    \item \textbf{Trigger}: Problemi di connessione o indisponibilità del database remoto.
\end{itemize}

\begin{comment}
% UC43 - Creazione collegamento tra note nel database
\subsubsection{UC 43: (Opzionale) Creazione collegamento tra note nel database}

% use case diagram
\begin{figure}[H]
    \centering
    \includegraphics[width=0.6\textwidth]{Collegamento_note_database.png}
    \caption{UC43 – Creazione collegamento tra note nel database}
\end{figure}

\begin{itemize}
    \item \textbf{Attore Principale}: Utente
    \item \textbf{Attore Secondario}: Database
    \item \textbf{Scenario Principale}: 
    \begin{enumerate}
        \item L'utente inserisce un riferimento a una nota esistente nel database;
        \item Il sistema crea una connessione logica tra i documenti nel database.
    \end{enumerate}
    \item \textbf{Precondizioni}: L'utente è autenticato e le note sono presenti sul database.
    \item \textbf{Postcondizioni}: Viene creato un collegamento bidirezionale tra le note.
    \item \textbf{Trigger}: L'utente vuole connettere logicamente due documenti remoti.
\end{itemize}

% UC43.1 - Visualizzazione errore sincronizzazione Database
\subsubsection{UC 43.1: Visualizzazione errore sincronizzazione Database}

% use case diagram
\begin{figure}[H]
    \centering
    \includegraphics[width=0.6\textwidth]{Errore_sincronizzazione_database_collegamento.png}
    \caption{UC43.1 – Visualizzazione errore sincronizzazione Database}
\end{figure}

\begin{itemize}
    \item \textbf{Attore Principale}: Utente
    \item \textbf{Scenario Principale}: 
    \begin{enumerate}
        \item L'utente tenta di collegare due note;
        \item Il sistema non riesce a scrivere il nuovo collegamento sul database;
        \item Il sistema mostra un messaggio di errore di sincronizzazione.
    \end{enumerate}
    \item \textbf{Precondizioni}: Procedura di collegamento avviata.
    \item \textbf{Postcondizioni}: Il collegamento non viene creato.
    \item \textbf{Trigger}: Problemi di integrità del database o di rete durante il collegamento.
\end{itemize}
\end{comment}